% concatenated from nakano_and_related.tex
\documentclass{amsart}
% \usepackage{hyperref,cleveref,amssymb,amsmath,color}
\usepackage{hyperref, cleveref, amsthm, amssymb, color, mathtools, thmtools}
%\usepackage{showkeys}


\usepackage{comment}
\usepackage{enumerate}

\usepackage{tikz, tcolorbox}
\usetikzlibrary{cd}

\input xy
\xyoption{all}

% \textheight 8.5in
% \textwidth 6 in
% \hoffset -0.5in
% %\oddsidemargin 0.25in
% \topmargin 0in

\theoremstyle{plain}
\newtheorem{theorem}{Theorem}[section]
\newtheorem{proposition}[theorem]{Proposition}
\newtheorem{corollary}[theorem]{Corollary}
\newtheorem{lemma}[theorem]{Lemma}
\newtheorem{conjecture}[theorem]{Conjecture}

\theoremstyle{definition}
\newtheorem{remark}[theorem]{Remark}
\newtheorem{question}[theorem]{Question}
\newtheorem{numbered}[theorem]{}
\newtheorem{example}[theorem]{Example}
\newtheorem{definition}[theorem]{Definition}

\numberwithin{equation}{section}

\crefname{theorem}{Theorem}{Theorems}
\crefname{proposition}{Proposition}{Propositions}
\crefname{section}{Section}{Sections}

% \renewcommand{\theenumi}{\roman{enumi}} % Use small romans for items
% \renewcommand{\labelenumi}{(\theenumi)}

\newcommand{\N}{\mathbb{N}}
\newcommand{\Z}{\mathbb{Z}}
\newcommand{\Q}{\mathbb{Q}}
\newcommand{\R}{\mathbb{R}}
\newcommand{\C}{\mathbb{C}}

\newcommand{\dist}{\textrm{dist}\,}
\newcommand{\spn}{\textrm{span}\,}
\newcommand{\conv}{\textrm{conv}\,}
\newcommand{\ran}{\textrm{ran}\,}
\newcommand{\rank}{\textrm{rank}\,}
\newcommand{\vr}{\varepsilon}
\newcommand{\ball}{\mathrm{B}}
\newcommand{\sphere}{\textrm{S}}
\newcommand{\Ave}{\textrm{Ave}}
\newcommand{\codim}{\mathrm{codim} \, }

\newcommand{\abs}[1]{\lvert#1\rvert}
\newcommand{\norm}[1]{\lVert#1\rVert}
\newcommand{\bigabs}[1]{\bigl\lvert#1\bigr\rvert}
\newcommand{\bignorm}[1]{\bigl\lVert#1\bigr\rVert}
\newcommand{\Bigabs}[1]{\Bigl\lvert#1\Bigr\rvert}
\newcommand{\biggabs}[1]{\biggl\lvert#1\biggr\rvert}
\newcommand{\Bignorm}[1]{\Bigl\lVert#1\Bigr\rVert}
\newcommand{\biggnorm}[1]{\biggl\lVert#1\biggr\rVert}
\newcommand{\Biggnorm}[1]{\Biggl\lVert#1\Biggr\rVert}
\newcommand{\term}[1]{{\textit{\textbf{#1}}}}   % To introduce a term
\newcommand{\marg}[1]{\marginpar{\tiny #1}}     % A margin note
\newcommand{\timur}[1]{{\textcolor{blue}{{\bf T.O:} #1}}}
\newcommand{\ana}[1]{{\textcolor{blue}{{\bf A.I:} #1}}}

\renewcommand{\mid}{\::\:}

\newcommand{\id}{{\textrm{id}}}

\DeclareSymbolFont{bbold}{U}{bbold}{m}{n}
\DeclareSymbolFontAlphabet{\mathbbold}{bbold}
\def\one{\mathbbold{1}}
%\def\one{\mathbb 1}

\DeclareMathOperator{\Range}{Range}
\DeclareMathOperator{\Sol}{Sol}
\DeclareMathOperator{\Span}{span}
\DeclareMathOperator{\Int}{Int}
\DeclareMathOperator{\FVL}{FVL}
\DeclareMathOperator{\FBL}{FBL}
\DeclareMathOperator{\FBLp}{FBL^{\mathit{(p)}}}

\newcommand{\Hp}{H_{w^*}^{p}}
\newcommand{\nip}{I_{w^*}^{p}}
\newcommand{\Ip}{\overline{I}_{w^*}^{p}}

\renewcommand{\le}{\leqslant}
\renewcommand{\ge}{\geqslant}

\renewcommand{\baselinestretch}{1.2}

\begin{document}

\title[Upper bound properties via operators]{Upper bound properties of free and related Banach lattices via operators}

\author{Anastasiia Ianina and Timur Oikhberg}

% \author[Oikhberg]{T.~Oikhberg} 
% \email{oikhberg@illinois.edu}
% \address{Dept. of Mathematics, University of Illinois, Urbana-Champaign, }

\address{Department of Mathematics, University of Illinois at Urbana-Champaign, Urbana, IL}
 \email{aianina2@illinois.edu}
 \email{oikhberg@illinois.edu}

\thanks{This work was partially supported by UIUC CRB award RB26177. The authors are grateful to M.~Taylor, P.~Tradacete, and V.~Troitsky for many stimulating conversations.}
%\keywords{Banach lattice, basic sequence, order convergence}
%\subjclass[2010]{Primary: 46B42. Secondary: 46B15, 46A40}
% 46B42 Banach lattices
% 46A40 Ordered topological linear spaces, vector lattices
% 46B15 Summability and bases

% \date{\today}

\begin{abstract}
It is known that the free $p$-convex Banach lattice on a Banach space $X$ can be represented as a space of functions on the unit ball of $X^*$.
In this way, it gives rise to certain related (larger) lattices. To investigate such lattices, we introduce a new tool, related to operators into $X$.
This tool is then used to (i) determine whether the lattices in question possess, or fail, properties involving upper bounds of upward directed sets -- namely, the Fatou property, and the related property of monotonic boundedness; (ii) investigate the regularity of embeddings between spaces in question.
We find connections between the Fatou-like properties of $\FBLp[X]$ and the Radon-Nikodym property of $X$.
In addition, we give an example of $X \subset Y$ such that $\FBLp[X]$ is not a regular sublattice of $\FBLp[Y]$.
% We attempt to determine under which conditions on a Banach space $X$ the corresponding free Banach lattice $\FBLp[X]$ possesses, or fails, the Fatou property, or the related property of monotonic boundedness. The answer turns out to be connected to the geometry of $X$, in particular to the Radon-Nikodym Property.
% The techniques we develop allow us to give an example of an isometric embedding $i : X \to Y$ so that the corresponding injection $\overline{i} : \FBLp[X] \to \FBLp[Y]$ is not order continuous. We also examine the spaces of functions $\Ip[X]$ and $\Hp[X]$, related to the functional representation of $\FBLp[X]$.
\end{abstract}

\maketitle

\section{Introduction and preliminaries}\label{sec:intro}

The theory of free Banach lattices has been advancing rapidly in the past several years. A free $p$-convex Banach lattice $\FBLp[X]$, over a Banach space $X$, is the (necessarily unique) Banach lattice $L$ such that (i) there exists an isometric embedding $\phi_{X,p}: X \to L$, and (ii) every bounded linear operator $T : X \to Z$, where $Z$ is a Banach lattice with $p$-convexity constant $1$, extends to a lattice homomorphism $\widehat{T}: L \to Z$, with $\|\widehat{T}\| = \|T\|$, as illustrated by the diagram below:
$$
\xymatrix{
 \FBLp[X] \ar@{-->}[dr]^{\widehat{T}}  &   \\
  X \ar[r]^{T} \ar@{^{(}->}[u]^{\phi_{X,p}}  & Z } 
$$
It is known (see, e.g.~\cite{OiTTT} or \cite{GSLTT}) that $\FBLp[X]$ admits the following functional representation. Consider $C(\ball(X^*))$, the space of weak$^*$-continuous functions on the closed unit ball $\ball(X^*)$.
For $x \in X$, define $\phi_{X,p}(x) = \delta_x : x^* \mapsto \langle x^*,x\rangle$.
The \emph{free vector lattice} on $X$ ($\FVL[X]$) is the vector sublattice of $C(\ball(X^*))$ generated by $\phi_{X,p}(X)$.
For $f \in C(\ball(X^*))$, define the ``$p$-summing'' norm $\|f\|_p$: % , equal to the infimum of all $C > 0$ so that the inequality $(\sum_{i=1}^n |f(x_i^*)|^p)^{1/p} \leq C \|(x_i^*)_{i=1}^n\|_{p, {\textrm{weak}}}$. 
\begin{equation}
\|f\|_p = \inf \big\{ C : \forall x_1^*, \ldots, x_n^* \in X^*, (\sum_{i=1}^n |f(x_i^*)|^p)^{1/p} \leq C \|(x_i^*)_{i=1}^n\|_{p, {\textrm{weak}}} \big\} .
\label{def of p-norm}
\end{equation}
Throughout this paper, for $z_1, \ldots, z_n \in Z$ we write $\|(z_i)_{i=1}^n\|_{p, {\textrm{weak}}} = \sup \big\{ \| \sum_i \alpha_i z_i \| : \|(\alpha_i)_i\|_{\ell_{p'}} \leq 1 \big\}$, and $1/p + 1/p' = 1$. Then $\FBLp[X]$ is the closure of $\FVL[X]$ in the norm $\| \cdot \|_p$. % we shall simply write $\| \cdot \|$ if confusion is unlikely. 

This functional representation leads naturally to the following larger lattices: 
\begin{itemize}
 \item $H_{w^*}[X]$ is the linear space of positively homogeneous functions $X^* \to \R$, weak$^*$-continuous on bounded sets;
 \item $\Hp[X]$ is the space of all $f \in H_{w^*}[X]$ with $\|f\|_p < \infty$;
 \item $I_{w^*}[X]$ (respectively, $\nip[X]$) is the ideal generated in $H_{w^*}[X]$ by $(\delta_x)_{x \in X}$ (respectively, by $\FBLp[X]$);
 \item $\Ip[X]$ is the closure of $I_{w^*}[X]$ (equivalently, of $\nip[X]$) in $\Hp[X]$.
\end{itemize}
When $p = \infty$, \cite[Proposition 2.2]{OiTTT} implies that the spaces $H_{w^*}^\infty[X]$, $\overline{I}_{w^*}^\infty[X]$, and $\FBL^{(\infty)}[X]$ all coincide with the space of positively homogeneous weak$^*$-continuous functions on $\ball(X^*)$, and $\| \cdot \|_\infty$ equals the $\sup$ norm. In general, $\FBLp[X] \subset \nip[X] \subset \Ip[X] \subset \Hp[X]$, and we shall see later that all the inclusions may be proper. See \cite{OiTTT} for a thorough study of free Banach lattices, and \cite{LaTr} for more information on the related function spaces introduced above.


The question of linking properties of a Banach space $X$ with those of the free Banach lattice built on it was raised in \cite{OiTTT}, where several such connections were established. The present work can be viewed as a step in this direction. In particular, we relate certain order-theoretic properties of free Banach lattices to the Radon-Nikodym property (RNP) of an underlying Banach space. We also study natural embeddings between free Banach lattices, continuing the investigations initiated in \cite{OiTTT} and \cite{AD}. Our approach to both problems is based on operators with range in $X$, as their adjoints naturally give rise to positively homogeneous functions on $\ball(X^*)$.

We focus on two order-theoretic properties of Banach lattices: the Fatou property and monotonic boundedness.
We say that a Banach lattice $Z$ has the \emph{$\lambda$-Fatou property} if for every upward directed set ${\mathcal{S}} \subset Z_+$ with supremum $z_0$, one has $\|z_0\| \leq \lambda \sup_{z \in {\mathcal{S}}} \|z\|$. When $\lambda = 1$, we simply say that $Z$ has the \emph{Fatou property}.

Following e.g.~\cite[Definition 4.2]{DLOT} or \cite[Section 10]{Taylor}, we say that a Banach lattice $Z$ is \emph{monotonically bounded} if every norm-bounded increasing net in $Z_+$ has an upper bound. Below we show that, in this case, there exists $\lambda \geq 1$ so that every norm-bounded increasing net $(z_\alpha) \subset Z_+$ has an upper bound $z_0$ with $\|z_0\| \leq \lambda \sup_\alpha \|z_\alpha\|$.
We shall then say that $Z$ is \emph{$\lambda$-monotonically bounded}. Clearly, $\lambda$-monotonic boundedness implies the $\lambda$-Fatou property.

If, instead of upward directed nets, we restrict our attention to increasing sequences, we obtain the definitions of \emph{countable $\lambda$-Fatou} etc.~properties.
\begin{remark}
 In \cite{ABT}, the term ``$\lambda$-strong Nakano'' is used to refer to monotonic boundedness. The term ``boundedly order bounded'' is also used sometimes.
\end{remark}

As for embeddings of lattices, we are especially concerned with their order continuity. Recall that for vector lattices $E$ and $F$ and an injective lattice homomorphism $j: E \to F$, $j(E)$ is called a \emph{regular sublattice} of $F$ if, for every set $A \subset E$ with supremum $a$ in $E$, $ja$ is the supremum of $j(A)$ in $F$. This is easily seen to be equivalent to $j$ being order continuous.
Specializing to free and related lattices, we consider certain ``natural" embeddings. Namely, for Banach spaces $X, Y$ and a bounded linear operator $i: X \to Y$, we consider $\overline{i}: \Hp[X] \to \Hp[Y]$ defined by $[\overline{i} f](y^*) = f(i^* y^*)$. %It is straightforward to verify that $\overline{i} f$ indeed belongs to $\Hp[Y]$). 
The argument of \cite{OiTTT-cor} shows that $i$ is injective if and only if $\overline{i}$ is. Note that, by \cite{OiTTT}, $\overline{i}(\FBLp[X]) \subset \FBLp[Y]$; consequently, $\overline{i}(\Ip[X]) \subset \Ip[Y]$.

The paper is organized as follows. After gathering some essential facts about ``operator-generated'' functions in \Cref{general} and examining monotonic boundedness in \Cref{m-b}, we show in \Cref{+nakano} that if $X$ is a dual space with the RNP and $1 \leq p < \infty$, then $\FBLp[X]$, $\Ip[X]$, and $\Hp[X]$ are 1-monotonically bounded. Under the same assumptions, $\Hp[X]=\Ip[X] = \nip[X]$.

In \Cref{operators:general}, we develop a general operator-based approach to detecting failure of the Fatou and monotonic boundedness properties in the free and related lattices. This approach also turns out to be useful in distinguishing between these spaces: we provide a sufficient condition for the inclusions $\nip[X] \subset \Ip[X] \subset \Hp[X]$ to be proper.

In the following sections, we apply our operator approach. In particular, in \Cref{no compactness}, we show that, for $1 \leq p < \infty$, under certain assumptions on $X$ stronger than the failure of the RNP, $\FBLp[X]$, $\Ip[X]$, and $\Hp[X]$ fail the $\lambda$-Fatou property for every $\lambda$.  The case $p=\infty$ is different: $\FBL^{(\infty)}[X]$ fails the $\lambda$-Fatou property unless $X$ is finite-dimensional. We further show that, for a dual space $X$ and $p \in [1,\infty)$, $\Hp[X]$ is 1-monotonically bounded if and only if $\Hp[X]$ has the Fatou property, which in turn is equivalent to $X$ having the wRNP (the weak Radon-Nikodym Property, described in \cite[Section 7.4]{Bour}). A similar characterization holds for $\Ip[X]$: if $X$ is a dual space, then $\Ip[X]$ has the Fatou property if and only if $X$ has the wRNP. However, in \Cref{james tree}, we show that there exists a dual space $X$ with the wRNP for which $\Ip[X]$ is not monotonically bounded and $\FBLp[X]$ fails the $\lambda$-Fatou property for every $\lambda$.

In \Cref{Ip into Hp}, we show that, for certain spaces $X$ failing the RNP, the inclusions $\nip[X] \subset \Ip[X] \subset \Hp[X]$ are proper. In contrast, \Cref{sup-norm} establishes that, for every Banach space $X$, the $\|\cdot\|_p$-unit balls of $I_{w^*}[X]$ and $\Ip[X]$ are $\|\cdot\|_{\infty}$-dense in the $\|\cdot\|_p$-unit ball of $\Hp[X]$. 

Finally, in \Cref{FBLp into Ip}, we investigate the regularity of natural embeddings between the lattices under consideration. We prove that $\FBLp[X]$ is not a regular sublattice of $\nip[X]$ unless $X$ is finite-dimensional. This sharpens a result of \cite{LaTr}, where additional assumptions on $X$ were imposed to show that $\FBLp[X]$ is a proper subspace of $\Ip[X]$. We also consider regularity in the context of embeddings between free Banach lattices. Given an isometric embedding $i : X \to Y$,  the question of whether $\overline{i}(\FBLp[X])$ must be a regular sublattice of $\FBLp[Y]$ has hitherto been open (for $p=\infty$, an affirmative answer was obtained in \cite{AD}; see also \cite{OiTTT-cor}). We answer this question in the negative.

Throughout, we rely on standard notation and facts. In particular, we denote by $\ball(X)$ and $\sphere(X)$ the closed unit ball and the unit sphere of $X$, respectively, and by $B(X, Y)$ the space of bounded linear operators from $X$ to $Y$. The required background on Banach spaces can be found, for instance, in \cite{AK}. For information about Banach lattices, we refer the reader to e.g.~\cite{M-N}, while the RNP is handled in e.g.~\cite{Bour}, \cite{DU}, \cite{DU-survey}, or \cite[Chapter 5]{BL}. % An outline of ${\mathcal L}_p$ spaces can be found in \cite[Section 23]{DeFl}.

% All the measures are assumed to be $\sigma$-additive, but not necessarily $\sigma$-finite. % We shall denote by $\ball(X)$ the closed unit ball of $X$. 
% Finally, we shall often refer to the \emph{$p$-summing}, \emph{$p$-integral}, and \emph{$p$-nuclear} norms of operators, denoted by $\pi_p( \cdot )$, $\iota_p( \cdot )$, and $\nu_p( \cdot )$, respectively; see \cite{DJT} for more information.


\section{Free Banach lattices and $p$-summing operators}\label{general}

Our exposition relies heavily on the ideals $\Pi_p$, $I_p$, and $N_p$ of \emph{$p$-summing}, \emph{$p$-integral}, and \emph{$p$-nuclear} operators; the corresponding norms are denoted by $\pi_p( \cdot )$, $\iota_p( \cdot )$, and $\nu_p( \cdot )$, respectively; see \cite{DJT} for more information.

We shall often use a special class of functions: for $T : X^* \to Z$, we define $\phi_T(x^*) = \|T x^*\|$. By \eqref{def of p-norm}, $\|\phi_T\|_p = \pi_p(T)$. The following result can be obtained by combining \cite[Section 7]{JLTTT} (see also \cite[Proposition 2.6]{OiGEO}) with Pietsch Factorization Theorem.

\begin{proposition}\label{FBLp norm}
For $p \in [1,\infty)$, and $\phi \in H_{w^*}[X]$, the following statements are equivalent:
\begin{enumerate}
\item $\phi \in \Hp[X]$, with $\| \phi \|_p \leq 1$.
\item There exist a probability measure $\mu$ and an operator $T : X^* \to L_p(\mu)$ so that $\iota_p(T) \leq 1$, and $|\phi| \leq \phi_T$.
\item There exist a Banach space $Y$ and an operator $T : X^* \to Y$ so that $\pi_p(T) \leq 1$, and $|\phi| \leq \phi_T$.
\end{enumerate}
    %Then $\phi \in \Hp[X]$, with $\| \phi \| \leq 1$, if and only if  there exist a probability measure $\mu$ and an operator $T : X^* \to L_p(\mu)$ so that $\iota_p(T) \leq 1$, and $|\phi| \leq \phi_T$. % $|\phi(x^*)| \leq \|Tx^*\|$ for every $x^* \in X^*$.
\end{proposition}

%The triangle inequality gives us:

%\begin{lemma}\label{difference}
%	For $T_1, T_2 : X^* \to Z$, $\phi_{T_1-T_2} \geq |\phi_{T_1} - \phi_{T_2}|$.
%\end{lemma}

Another simple observation:

\begin{lemma}\label{Cauchy seq}
	If $(f_i)$ is a Cauchy sequence in $\FBLp[X]$, converging pointwise to a function $f$, then $f$ belongs to $\FBLp[X]$, and is the limit of $(f_i)$ in the norm of the latter space.
\end{lemma}

The following lemma will be used repeatedly throughout the text: 

\begin{lemma}\label{weak* continuity}
For $S \in B(Z,Y)$, the following are equivalent:
\begin{enumerate}
   \item $S$ is compact.
   \item $S^*$ is weak$^*$ to norm continuous on bounded subsets of $Y^*$.
   \item The function $F : \ball(Y^*) \to \R : y^* \mapsto \|S^* y^*\|$ is weak$^*$ continuous at $0$.
\end{enumerate}
% If a map $S : Z \to Y$ is compact, then $S^*$ is weak$^*$ to norm continuous on bounded subsets of $Y^*$.
\end{lemma}

\begin{proof}
$(1) \Rightarrow (2)$:
Let $q : \ell_1(I) \to Z$ be a quotient map, where $I$ is a sufficiently large index set. The space $\ell_\infty(I) = \ell_1(I)^*$ has the (Metric) Approximation Property, hence, by \cite[Theorem 1.e.5]{LT1}, there exists a net of finite rank maps $T_u : \ell_1(I) \to Y$, converging to $Sq$ in norm. Clearly $T_u^*$ is weak$^*$ to norm continuous, hence so is $q^* S^*$. Now recall that, for $y^* \in Y^*$, $\|S^* y^*\| = \|q^* S^* y^*\|$. 

$(2) \Rightarrow (3)$ is straightforward.

$(3) \Rightarrow (1)$: If $S$ is as in (3), then for any $\vr > 0$ there exist $y_1, \ldots, y_n \in Y$  so that $\|S^* y^*\| < \vr$ whenever $y^* \in \ball(Y^*)$ satisfies $\vee_i |\langle y^*, y_i \rangle| < 1$.
Then in particular, $\|S^*|_{E^\perp}\| \leq \vr$, where $E = \spn[y_1, \ldots, y_n]$. The former quantity equals $\|S^* q_E^*\| \leq \vr$; here $q_E : Y \to Y/E$ is the quotient map, and $q_E^* : E^\perp \to Y^*$ is the canonical embedding.
So for every $\vr > 0$ there exists a finite dimensional $E \subset Y$ so that $\|q_E S\| \leq \vr$. This implies that $S$ is compact.
\end{proof}

Combining the above with the definition of the norm in $\Hp[X]$, we obtain:

\begin{corollary}\label{compact_adjoint}
    If $T \in B(Y, X)$ is compact and $T^* \in \Pi_p(X^*, Y^*)$, then $\phi_{T^*} \in \Hp[X]$ and $\|\phi_{T^*}\|_p = \pi_p(T^*)$.
\end{corollary}

For finite rank operators, more can be said:

\begin{lemma}\label{fin rank ops}
Suppose $X$ is a Banach space, $E$ is finite dimensional, and $u \in B(X^*,E)$ is weak$^*$ continuous. Then $\phi_u \in \FBLp[X]$, and $\|\phi_u\| = \pi_p(u)$. Moreover, for every $\delta > 0$ there exists $\psi \in \FVL[X]_+$ so that $(1-\delta) \psi \leq (1-\delta) \phi_u \leq \psi$.
\end{lemma}

\begin{proof}
By scaling, we can assume that $\pi_p(u) = 1$. It is straightforward that $\phi_u$ is weak$^*$ continuous on bounded sets, and $\|\phi_u\|_{p} = \pi_p(u) = 1$. It remains to show that $\phi_u$ lies in $\FBLp[X]$ -- that is, for any $\vr > 0$ there exists $\psi \in \FVL[X]$ so that $\|\phi_u - \psi\| < \vr$.
To this end, find $N$ and a contraction $v \in B(E,\ell_\infty^N)$, so that $\|ve\| \geq 1 - \delta = 1 - \vr/2$ for any $e \in E$ of unit norm. Write $ve = (\langle e_j^*,e \rangle)_{j=1}^N$. Let $\psi(x^*) = \|vux^*\|$, then clearly $\psi = \vee_{j=1}^N \big| \delta_{u^* e_j^*} \big|$.
The weak$^*$-continuity of $u$ implies $u^* e_j^* \in X$, so $\psi \in \FVL[X]$. Also, $\psi \leq \phi_u \leq (1-\delta)^{-1} \psi$, hence $0 \leq \phi_u - \psi \leq \delta(1-\delta)^{-1} \phi_u$, giving us the desired estimate for $\|\phi_u - \psi\|$.
\end{proof}

\begin{corollary}\label{nuclear in FBL}
    If $u: X^* \to Y$ is approximable in the $p$-summing norm by a sequence of weak$^*$ continuous finite-rank operators $u_n : X^* \to Y$, then $\phi_u \in \FBLp[X]$ with $\|\phi_u\| = \pi_p(u)$.
\end{corollary}

We shall see later that the conclusion $\phi_u \in \FBLp[X]$ may fail if $u$ is merely approximable (in the operator norm).

\section{Introduction to monotonic boundedness}\label{m-b}

We gather the essential facts about monotonically bounded Banach lattices. For applications of this notion to regular operators, see \cite[Section 4]{DLOT}, and for other related properties, see \cite{Taylor}.

\begin{proposition}
 If a Banach lattice $X$ is monotonically bounded (respectively, countably monotonically bounded), then it is $\lambda$-monotonically bounded (respectively, countably $\lambda$-monotonically bounded) for some $\lambda$.
\end{proposition}

\begin{proof}
 We prove the ``monotonically bounded'' case, the countable setting is handled similarly.
 
 Suppose, for the sake of contradiction, that $X$ is monotonically bounded, but not $\lambda$-monotonically bounded for any $\lambda$. Then for every $n$ these exists an increasing net $(z^{(n)}_\alpha)_{\alpha \in {\mathcal{A}}_n}$ so that $\sup_{\alpha \in {\mathcal{A}}_n} \|z^{(n)}_\alpha\| < 2^{-n}$, yet $\|z_0\| > 2^n$ whenever $z_0$ is an upper bound for this net.
 
 Let $\mathcal{A}$ the set of all finite sequences $(\alpha_1, \ldots, \alpha_m)$, with $m \in \N$ and $\alpha_i \in {\mathcal{A}}_i$. Equip $\mathcal{A}$ with the ``lexicographic'' order: $(\alpha_1, \ldots, \alpha_m) \prec (\beta_1, \ldots, \beta_n)$ if $m \leq n$, and $\alpha_i \leq \beta_i$ for $1 \leq i \leq m$. Then the net $z_{(\alpha_1, \ldots, \alpha_m)} = \sum_{i=1}^m z^{(i)}_{\alpha_i}$ is upward directed, and norms of its elements are bounded by $1$. On the other hand, it cannot have an upper bound.
\end{proof}

We can also observe a connection between monotonic boundedness and other related properties.

\begin{proposition}\label{connections}
 \begin{enumerate}
  \item If a Banach lattice $X$ is $\lambda$-monotonically bounded, then it has the $\lambda$-Fatou property. The same holds for the ``countable'' versions of these properties.
  \item If $X$ is monotonically complete, then it is monotonically bounded.
 \end{enumerate}
The converse statements are false.
\end{proposition}

Recall that $X$ is called \emph{monotonically complete} if every norm-bounded increasing net in $X$ has a supremum. For properties of such lattices, see \cite[Section 2.4]{M-N}.

\begin{proof}
 The forward statements are straightforward. It remains to provide counterexamples.
 
 (1) Let $X$ be the space of continuous functions on $[0,1]$ which vanish at $0$. By \cite{Wick07}, $X$ has the Fatou property. However, it is easy to see it isn't monotonically bounded.
 
 (2) It is easy to see that $X = C[0,1]$ works. 
\end{proof}

By \cite[Theorem 2.4.19]{M-N}, any dual of a Banach lattice is monotonically complete and has the Fatou property; consequently, it is also 1-monotonically bounded.
On the other end of the spectrum, an AM-space is 1-monotonically bounded if and only if it is lattice isomorphic to a $C(K)$ space, for some Hausdorff compact $K$ \cite{ABT}.


\begin{comment}

\begin{proposition}
 An AM-space is 1-monotonically bounded if and only if it is lattice isomorphic to a $C(K)$ space, for some Hausdorff compact $K$.
\end{proposition}

This is actually proved in \cite{ABT}, we present an alternative approach.
\marg{\ana{Check if the appeal to Nakano is necessary here}}
\begin{proof}
 Clearly, $C(K)$ is 1-monotonically bounded. Conversely, suppose an AM-space $X$ is 1-monotonically bounded. By \Cref{connections}, $X$ has the strong Nakano property, hence, by \cite{Wick07}, it is lattice isometric to a $C_0(\Omega)$ space, for some locally compact $\Omega$. 
 If $\Omega$ is not compact, consider the net ${\mathcal{K}}$ of pairs of compact subsets of $\Omega$, $\widetilde{K} = (K_1, K_2)$, so that $K_1$ is contained in the interior of $K_2$ (by local compactness, for any $K_1$ there exists a $K_2$ like this).
 Define order on $\mathcal K$: $(K_1, K_2) \prec (L_1,L_2)$ if $K_2 \subset L_1$.
 For each $\widetilde{K} \in {\mathcal K}$, use Urysohn Lemma to find a continuous function $f_{\widetilde{K}} : \Omega \to [0,1]$, which is equal to $1$ on $K_1$, and vanishes outside of $K_2$.
 Clearly $(f_{\widetilde{K}})_{\widetilde{K} \in {\mathcal{K}}}$ is a norm bounded increasing net, with no upper bound in $C_0(\Omega)$.
\end{proof}

\end{comment}

Finally, we examine monotonic boundedness in separable lattices. The proof of the following proposition closely imitates that of \cite[Lemma 2.5]{ABT}.

\begin{proposition}\label{sep mon bound}
 A separable Banach lattice is $\lambda$-monotonically bounded if and only if it is countably $\lambda$-monotonically bounded.
\end{proposition}

We should also mention \cite[Proposition 12.3]{dPWi}: if a monotonically bounded Banach lattice $Y$ is embedded into a Banach lattice $X$ as an ideal, then $Y$ is a projection band in $X$.


\section{A class of monotonically bounded free lattices}\label{+nakano}

Throughout this section, we assume that $1 \leq p < \infty$. Our first goal is to determine under which assumptions $\Hp[X]$ is monotonically bounded. As the next proposition shows, this turns out to be closely related to the weak Radon-Nikodym property of an underlying Banach space $X$ (see \cite[Section 7.4]{Bour} for background on the wRNP).
\begin{theorem}\label{wRNP_Hp}
 If $X$ is a dual Banach space with the wRNP, then $\Hp[X]$ is 1-monotonically bounded.
\end{theorem}
In fact, in \Cref{duals and Hp}, we shall show that, for a dual Banach space $X$, $\Hp[X]$ is monotonically bounded if and only if X has wRNP (in which case, it is 1-monotonically bounded).

We shall make use of ultraproducts (see \cite{Hein} for an introduction into the topic).

\begin{lemma}\label{ultraproduct}
Suppose $\mathfrak U$ is an ultrafilter on a set $I$, $X_i, Y_i$ $(i \in I)$ are Banach spaces, and $T_i : X_i \to Y_i$ are linear operators, so that $\sup_{i \in I} \pi_p(T_i) < \infty$. Then $\pi_p(\prod_{\mathfrak U} T_i) \leq \lim_{\mathfrak U} \pi_p(T_i)$.
\end{lemma}

\begin{remark}
 The inequality in \Cref{ultraproduct} may be strict. For instance, take $p=2$, $I = \N$, $X_i = Y_i = \ell_2$ for any $i$, and suppose $\mathfrak U$ is a free ultrafilter on $\N$. Let $T_i = i^{-1/2} P_i$, where $P_i$ is an orthogonal projection of rank $i$. We know that $\pi_2(T_i) = 1$ for any $i$, yet $\prod_{\mathfrak U} T_i = 0$.
\end{remark}


\begin{proof}
By scaling, we can assume that $\lim_{\mathfrak U} \pi_p(T_i) = 1$.
We have to show that, for any $\vr > 0$ and $(\overline{x}_j)_{j=1}^N \in \prod_{\mathfrak U} X_i$ with $\|(\overline{x}_j)\|_{p,\textrm{weak}} < 1$, we have $\sum_j \| (\prod_{\mathfrak{U}} T_i) \overline{x}_j\|^p \leq (1+\vr)^{2p}$.
Note that, if $\overline{x}_j$ is represented by $(x_{ji})_{i \in I}$ ($x_{ji} \in X_i$), then $(\prod_{\mathfrak{U}} T_i) \overline{x}_j$ is the equivalence class of $(T_i x_{ji})_{i \in I}$, hence it suffices to show the existence of $U \in {\mathfrak{U}}$ so that $\sum_j \|T_i x_{ji}\|^p \leq (1+\vr)^{2p}$ for any $i \in U$.

To this end, for each $j$ find $(x_{ji})_{i \in I} \in (\sum_{i \in I} X_i)_\infty$ so that $\overline{x}_j$ is the equivalence class of $(x_{ji})_{i \in I}$ (modulo the class $N_{\mathfrak{U}}$ of all $(y_i)$ with $\lim_{\mathfrak{U}} \|y_i\| = 0$). Let $C = \max_j \sup_i \|x_{ji}\|$.
Denote by $S$ the set of all $\widetilde{\alpha} = (\alpha_j)_{j=1}^N$ so that $\sum_j |\alpha_j|^q \leq 1$ ($1/p + 1/q = 1$). Let $S_0$ be a finite subset of $S$ with the property that for any $\widetilde{\alpha} = (\alpha_j)_{j=1}^N \in S$ there exists $\widetilde{\beta} = (\beta_j)_{j=1}^N \in S_0$ so that $\max_j |\alpha_j - \beta_j| < \vr/(4N\max\{C,1\})$.

For $\widetilde{\beta} = (\beta_j) \in S_0$ find $y_{\widetilde{\beta}} = (y_{\widetilde{\beta}i}) \in N_{\mathfrak{U}}$ so that $\|\sum_j \beta_j x_{ji} - y_{\widetilde{\beta}i}\| < 1 + \vr/4$ for any $i$. The set $I_{\widetilde{\beta}} = \{i \in I : \|y_{\widetilde{\beta}i}\| < \vr/4\}$ belongs to ${\mathfrak{U}}$, hence so does $I_0 = \cap_{\widetilde{\beta} \in S_0} I_{\widetilde{\beta}}$. 
By the triangle inequality, for any $i \in I_0$ and $\widetilde{\beta} \in S_0$, we have $\big\| \sum_j \beta_j x_{ji} \big\| \leq 1 + \vr/2$. A further application of the triangle inequality gives us $\big\| \sum_j \alpha_j x_{ji} \big\| \leq 1 + \vr$ for any $\widetilde{\alpha} = (\alpha_j) \in S$ and $i \in I_0$.
Further, there exists $I_1 \in {\mathfrak{U}}$ so that $\pi_p(T_i) < 1+\vr$ for $i \in I_1$. From this we conclude that, for any $i \in U = I_0 \cap I_1$, 
$$
\big( \|\sum_j \|T_i x_j\|^p \big)^{1/p} \leq \pi_p(T_i) \sup_{\widetilde{\alpha} \in S} \big\| \sum_j \alpha_j x_{ji} \big\| \leq (1+\vr)^2 ,
$$
and our task is complete.
\end{proof}

\begin{proof}[Proof of \Cref{wRNP_Hp}]
Given an increasing net $(f_i)_{i \in I} \subset \ball(\Hp[X])_+$, we shall find $f \in \ball(\Hp[X])_+$ so that $f \geq f_i$ for every $i$.
By \Cref{FBLp norm}, for each $i$, there exist a probability measure $\mu_i$ and an operator $T_i : X^* \to L_p(\mu_i)$ with $\iota_p(T_i) \leq 1$, and $\|T_i x^*\| \geq f_i(x^*)$ for any $x^* \in X^*$.

Consider the \emph{order filter} on $I$, which consists of sets $S \subset I$ for which there exists $i \in S$ so that $j \in S$ whenever $j \geq i$. An ultrafilter containing such a filter shall be called \emph{order consistent}.
% Note that any free ultrafilter on $\N$ is order consistent, as it must contain complements of finite sets.
Suppose ${\mathcal{U}}$ is an order consistent ultrafilter on $I$, then, by \Cref{ultraproduct}, $\prod_{\mathcal{U}} T_i : (X^*)^{\mathcal{U}} \to (L_p(\mu_i))^{\mathcal{U}}$ has $p$-summing norm not exceeding $\lim_{\mathcal U} \pi_p(T_i) \leq 1$. 
Let $S_0$ be the restriction of the aforementioned operator to $X^*$ (canonically embedded into its ultrapower), then $\pi_p(S_0) \leq 1$, and, for any $x^* \in X^*$, $\|S_0 x^*\| = \lim_{{\mathcal{U}}} \|T_i x^*\| \geq \vee_i f_i(x^*)$. 
By Pietsch Factorization Theorem, there exists an operator $S : X^* \to L_p(\mu)$ so that $\iota_p(S) \leq 1$, and $\|Sx^*\| \geq \vee_i f_i(x^*)$ for any $x^* \in X^*$.

Now suppose $X_*$ is a predual of $X$, and let $T = S|_{X_*}$. By \cite[Theorem II.1]{Car}, $T$ is compact, therefore, since $\pi_p(T^{**}) = \pi_p(T) \leq 1$, \Cref{compact_adjoint} yields that $\phi_{T^{**}} \in \Hp[X]$ and $\|\phi_{T^{**}}\|_p \leq 1$.

It remains to show that $\phi_{T^{**}} \geq f_i$ for every $i$. Indeed, any $x^* \in \ball(X^*)$ is the weak$^*$ limit of a net $(z_\alpha) \subset \ball(X_*)$ . As $T^{**}$ is weak$^*$ continuous, $\phi_{T^{**}}(x^*) = \lim_\alpha \phi_{T^{**}}(z_\alpha)$. Also, for any $i$, $f_i(x^*) = \lim_\alpha f_i(z_\alpha)$. Finally, for each $\alpha$, $\phi_{T^{**}}(z_\alpha) = \|S z_\alpha\| \geq f_i(z_\alpha)$, which leads to the desired conclusion.
\end{proof}

Before continuing our discussion of monotonic boundedness, we make a brief detour to investigate a question raised in \cite{LaTr}. Specifically, we prove the following.

\begin{theorem}\label{la-tr}
 If $X$ is a dual Banach space with the RNP, and $1 \leq p < \infty$, then for any $f \in \ball(\Hp[X])_+$ there exists $g \in \ball(\FBLp[X])_+$ such that $f \leq g$. As a consequence, $\Hp[X] = \nip[X]$.
\end{theorem}

In \Cref{Ip into Hp} we shall see examples of spaces $X$ with $\nip[X] \subsetneq \Ip[X] \subsetneq \Hp[X]$.

\begin{proof}
 Consider $f \in \ball(\Hp[X])_+$. Find a probability measure $\mu$ and an operator $S : X^* \to L_p(\mu)$ with $\iota_p(S) \leq 1$, and $\|S x^*\| \geq f(x^*)$ for any $x^* \in X^*$.
 
 Let $S_0 = S|_{X_*} : X_* \to L_p(\mu)$ (as before, $X_*$ is a predual of $X$), then $\pi_p(S_0) \leq 1$. By Pietsch Factorization Theorem, there exists an operator $T : X_* \to L_p(\mu)$ so that $\iota_p(T) \leq 1$, and $\|Tx_*\| \geq f(x_*)$ for any $x_* \in X_*$. In fact, in the terminology of \cite{Car}, $T$ is strictly $p$-integral, since $L_p(\mu)$ is contractively complemented in its second dual. 
 Therefore, by \cite[Theorem II.4]{Car}, $T$ is $p$-nuclear with $\nu_p(T) \leq 1$. Consequently, there exists a sequence of finite-rank operators $T_n: X_* \to L_p(\mu)$ such that $\nu_p(T - T_n) \to 0$. Note that $\pi_p(T^{**} - T_n^{**}) = \pi_p(T - T_n) \leq \nu_p(T - T_n)$, so by \Cref{nuclear in FBL}, $\phi_{T^{**}} \in \FBLp[X]$, and $\|\phi_{T^{**}}\|_p \leq 1$.
 
  Further, $f(x_*) \leq \|T x_*\| = \|T^{**} x_*\| = \phi_{T^{**}}(x_*)$ for $x_* \in X_*$. As both $f$ and $\phi_{T^{**}}$ are weak$^*$ continuous on bounded sets, Goldstine Theorem implies $f \leq \phi_{T^{**}}$, as desired.
\end{proof}

As the RNP implies  the wRNP, the following is an immediate consequence of \Cref{wRNP_Hp,la-tr}.
\begin{corollary}\label{dual-with-RNP}
 If $X$ is a dual Banach space with the RNP, then both $\FBLp[X]$ and $\Ip[X]$ are 1-monotonically bounded.
\end{corollary}

In particular, we have the following:

\begin{corollary}\label{good cases}
If $X$ is a separable dual space or $X$ is reflexive, then both $\FBLp[X]$ and $\Ip[X]$ are 1-monotonically bounded.
\end{corollary}

\begin{proof}
In both cases, $X$ has the RNP (see, for instance, \cite{DU} or \cite[Chapter 5]{BL}).
\end{proof}

\begin{remark}\label{lattice case}
By \cite[Corollary 5.4.21]{M-N}, a separable Banach lattice $X$ has the RNP iff it is a dual of a Banach lattice. Thus, if $X$ is a separable Banach lattice with the RNP (equivalently, with the wRNP, see \cite{GS}), then the lattices $\FBLp[X], \Ip[X], \Hp[X]$ are 1-monotonically bounded.
\end{remark}

\section{Functions on $X^*$ generated by operators}\label{operators:general}

In the preceding section, we have already used operators to investigate spaces of functions on $X^*$. We will pursue this further in this section, and below.
The main set-up will be as follows. Suppose $Z, X$ are Banach spaces, $T \in B(Z,X)$, and $Z_1 \subset Z_2 \subset \ldots$ are finite dimensional subspaces of $Z$. We say that $(Z_k)$ is a \emph{sequential finite dimensional resolution} (\emph{SFDR}) for $T$ if $\cup_k Z_k$ is $1$-norming for $T^*(X^*)$ -- that is, the equality
$$
\|T^* x^*\| = \sup_k \sup \big\{ \langle T^* x^*, z \rangle : z \in \ball(Z_k) \big\} 
$$
holds for any $x^* \in X^*$. Further, let $T_k = T \big|_{Z_k}$. A simple observation is in order:

\begin{lemma}\label{l:simple properties}
In the above notation:
\begin{enumerate}
    \item For every $x^* \in X^*$, $\|T_k^* x^*\| \nearrow \|T^* x^*\|$.
    \item $\pi_p(T_k^*) \nearrow \pi_p(T^*)$.
\end{enumerate}
\end{lemma}

In item (2), $\pi_p(T^*)$ need not be finite: throughout this paper, we set $\pi_p(S) = \infty$ if $S$ is not $p$-summing.

In a similar fashion, we say that $(Z_k)$ is a \emph{net finite dimensional resolution} (\emph{NFDR}) for $T$ if $(Z_k)_{k \in {\mathcal K}}$ is an increasing net of finite dimensional subspaces of $X$, and $\cup_k Z_k$ is $1$-norming for $T^*(X^*)$.
Results similar to the above holds for NFDR as well. Note that any operator has an NFDR. However, most operators we shall consider (in particular, those with separable domain) will actually possess SFDR.

\begin{comment}

\begin{remark}
We could have given similar definitions for an increasing net $(Z_\alpha)$ of finite dimensional subspaces of $Z$, instead of a sequence $(Z_k)$. 
We do not do so, however, since we have no applications for this more general set-up.
\end{remark}

    
\end{comment}

% Throughout, we shall consider the functions $\phi(x^*) = \|T^* x^*\|$ and $\phi_k(x^*) = \|T_k^* x^*\|$. From the above, $\phi \nearrow \phi_k$.

\begin{comment}

\begin{lemma}
Suppose $E$ is finite dimensional, and $S \in B(E,X)$. For $x^* \in X^*$ let $\psi(x^*) = \|S^* x^*\|$.
Then the function $\psi$ belongs to $\FBLp[X]$, for every $p$. Moreover, for any $C>1$ there exists $g \in \FVL[X]$ so that $g \leq \phi_k \leq Cg$.
\end{lemma}

\begin{proof}
Find a contraction $v : E^* \to \ell_\infty^N$ so that $\|vz\| \geq \|z\|/C$, for any $z \in E^*$. For $x^* \in X^*$ let $g(x^*) = \|v S^* x^*\|$. Clearly $g \leq \psi \leq C g$.
To show that $g \in \FVL[X]$, denote by $(f_j)$ the canonical basis of $\ell_1^N$. For any $x^* \in X^*$,
$g(x^*) = \vee_{j=1}^N |\langle vS^* x^*, f_j \rangle| = \vee_{j=1}^N |\langle x^*, S v^* f_j \rangle|$, which clearly represents a function from $\FVL[X]$.
\end{proof}
    
\end{comment}

\begin{theorem}\label{no compact no fatou}
Suppose there exists a Banach space $Z$ and a non-compact $T \in B(Z,X)$ such that $T^* \in \Pi_p(X^*,Z^*)$.
Then the lattices $\FBLp[X], \Ip[X], \Hp[X]$ fail the $\lambda$-Fatou property, for any $\lambda > 0$.
If, moreover, $T$ has an SFDR, then the above lattices fail the countable $\lambda$-Fatou property.
\end{theorem}

\begin{proof}
We consider the general case, as the one with the SFDR is similar.
Let $E$ be any of the three lattices in question. By scaling, we can assume $\pi_p(T^*) = 1$, and hence $\|T\| \leq 1$. Let $(Z_k)$ be an NFDR for $T$.
As before, write $T_k = T \big|_{Z_k}$, $\phi(x^*) = \|T^* x^*\|$, $\phi_k(x^*) = \|T_k^* x^*\|$. Note that, by \Cref{fin rank ops}, $\phi_k \in \FBLp[X]$ for all $k$.
Pick $\mu > \lambda$, and let $f_k = \phi_k \wedge \mu \big|\delta_{x_0}\big|$, where $x_0 \in X$ has unit norm.
We claim that $\mu \big|\delta_{x_0}\big| = \vee_k f_k$ in $E$. As $\|\mu \big|\delta_{x_0}\big|\| = \mu > \lambda$, and $\|f_k\| \leq \|\phi_k\| \leq \|\phi\| = 1$ for every $k$, we shall be done once the claim is established. Thus, we shall show that if $g \in E$ is such that $g \geq f_k$ for every $k$, then $g \geq \mu \big|\delta_{x_0}\big|$.
In other words, $g(x^*) \geq \mu |\langle x^*, x_0 \rangle|$ for every norm one $x^* \in X^*$.

By \Cref{weak* continuity}, $\phi|_{\ball(X^*)}$ is discontinuous at $0$, so there exists $c > 0$ and a weak$^*$ null net $(x_\alpha^*) \subset \ball(X^*)$ such that $\phi(x_\alpha^*) > c$ for each $\alpha$. Thus, for each $\alpha$, there exists $K(\alpha)$ such that $\phi_k(x_\alpha^*) > c$ for $k \geq K(\alpha)$.

Fix a norm one $x^* \in X^*$ and pick $C > (1 + \mu)/c$. Then
$$
\phi_k(x^* + C x_\alpha^*) = \|T_k^*(x^* + C x_\alpha^*)\| \geq C \|T_k^* x_\alpha^*\| - \|T_k^* x^*\| ,
$$
and hence $\phi_k(x^* + C x_\alpha^*) \geq Cc-1$ for $k \geq K(\alpha)$. Now fix $\sigma > 0$, and find $\alpha_0$ such that $\big| \langle x_\alpha^* , x_0 \rangle \big| < \sigma/C$ for $\alpha \geq \alpha_0$. For such $\alpha$,
$$
\mu \big| \delta_{x_0} \big| (x^* + C x_\alpha^*) = \mu \big| \langle x^* + C x_\alpha^* , x_0 \rangle \big| > \mu \big| \langle x^* , x_0 \rangle \big| - \sigma .
$$
By our choice of $C$, for $\alpha \geq \alpha_0$ and $k \geq k(\alpha)$,
$$
f_k (x^* + C x_\alpha^*) \geq (Cc - 1) \wedge \big( \mu \big| \langle x^* , x_0 \rangle \big| - \sigma \big) = \mu \big| \langle x^* , x_0 \rangle \big| - \sigma .
$$
Consequently, for $\alpha \geq \alpha_0$, $g (x^* + C x_\alpha^*) \geq \mu \big| \langle x^* , x_0 \rangle \big| - \sigma$. Passing to the limit along the net, we conclude that $g (x^*) \geq \mu \big| \langle x^* , x_0 \rangle \big| - \sigma$. It remains to recall that $\sigma$ is arbitrarily small.
\end{proof}

\begin{remark}
The same proof goes through with $\mu \big| \delta_{x_0} \big|$ replaced by an arbitrary $h \in \FBLp[X]_+$.
\end{remark}


To perform a finer analysis, suppose $Z, X$ are Banach spaces, and $X_0 \subset X$ is dense. We say that $T \in B(Z,X)$ is \emph{dual $p$-persistent relative to $X_0$} if there exists $c > 0$ so that, for any $x_1, \ldots, x_n \in X_0$, we have $\pi_p(T^*|_{\cap_i x_i^\perp}) > c$.
We shall call $T$ \emph{dual $p$-persistent} if it is dual $p$-persistent with respect to some dense set $X_0 \subset X$.

Note that $T$ is dual $p$-persistent relative to $X$ if $T^*$ is not $p$-summing: indeed, then the restriction of $T^*$ onto any subspace of finite codimension will fail to be $p$-summing.

\begin{comment}
for every finite dimensional $E \subset X$, $\pi_p(T^*|_{E^\perp}) > c$. Here, as before, we use the convention of $\pi_p(S) = \infty$ is $S$ is not $p$-summing.

The following result makes verifying the dual $p$-persistence easier:

\begin{lemma}\label{dense subset}
Suppose $X_0$ is a dense subset of $X$. Then $T$ is dual $p$-persistent if and only if there exists $c > 0$ such that, for any $x_1, \ldots, x_n \in X_0$, $\pi_p(T^*|_{\cap_i x_i^\perp}) > c$.
\end{lemma}

\begin{proof}
If $T^*$ is not $p$-summing, then its restriction on any subspace of finite codimension is not $p$-summing either. If $T^*$ is $p$-summing, then 

FILL IN
\end{proof}

\end{comment}

\begin{theorem}\label{failure of fatou}
Suppose there exists a Banach space $Z$ and a dual $p$-persistent $T \in B(Z,X)$ such that $T^* \in \Pi_p(X^*,Z^*)$. Then $\FBLp[X]$ fails the $\lambda$-Fatou property, for any $\lambda$, and $\Ip[X]$ is not monotonically bounded.
\end{theorem}

The proof relies on:


\begin{lemma}\label{lipschitz}
 If $P : \R^n \to \R$ is lattice polynomial, and $p \in [1,\infty)$, then there exists $K > 0$ so that, for any $t = (t_i)_{i=1}^n$ and $s = (s_i)_{i=1}^n$, we have $|P(t) - P(s)|^p \leq K \sum_i |t_i - s_i|^p$.
\end{lemma}

\begin{proof}
 By the standard H\"older Inequality, it suffices to show that for any $P$ there exists $K_0 = K_0(P)$ so that, for any $t,s$ as above, 
 \begin{equation}
  |P(t) - P(s)| \leq K_0 \sum_i |t_i - s_i| .
  \label{eq:lipschitz}
 \end{equation}
 This we establish by induction on the number of lattice operations involved in $P$.
 The base of induction reduces to considering $P(t) = t_i$, which clearly satisfies \eqref{eq:lipschitz}, with $K_0 = 1$.
 
 For the inductive step, it suffices to show that if $P_1, P_2$ satisfy \eqref{eq:lipschitz}, then so do $P_1 + P_2$, $P_1 \vee P_2$, and $a P_1$ for any $a \in \R$ (with a different constant $K_0$). This, however, is straightforward. For instance, for $P = P_1 \vee P_2$, note that $|P(t) - P(s)| \leq |P_1(t) - P_1(s)| + |P_2(t) - P_2(s)|$, hence $K_0(P) \leq K_0(P_1) + K_0(P_2)$.
\end{proof}

\Cref{lipschitz} directly implies:

\begin{corollary}\label{approximation}
Suppose $X_0$ is a dense subset of $X$. Denote by $\mathcal Y$ the family of all lattice expressions $F(\delta_{x_1}, \ldots, \delta_{x_n})$, with $x_1, \ldots, x_n \in X_0$. Then:
\begin{enumerate}
\item For any $\vr > 0$ and $f \in \FBLp[X]$ there exists $g \in \mathcal Y$ such that $\|f-g\| < \vr$. If $f$ is positive, then $g$ can be taken positive as well.
\item For any $\vr > 0$ and $f \in \Ip[X]_+$ there exist $g \in {\mathcal Y}_+$ and $h \in \Hp[X]_+$ such that $\|h\| < \vr$ and $f \leq g+h$.
\end{enumerate}
\end{corollary}

\begin{lemma}\label{go far to get nothing}
Suppose $X_0$ is dense in $X$. Then, for any $g \in \Ip[X]$ and $\vr > 0$, there exist $x_1, \ldots, x_m \in X_0$ so that, for every $x_1^*, \ldots, x_n^* \in \cap_{i=1}^m x_i^\perp$ with $\|(x_j^*)\|_{p,\textrm{weak}} \leq 1$, we have $\sum_j |g(x_j^*)|^p < \vr^p$.
\end{lemma}

\begin{proof}
We can and do assume $g \geq 0$. Fix $\vr > 0$, %and write $g = g_1 + g_2$ so that $\|g_2\| < \vr/2$, and $g_1 \in \nip[X]_+$ --that is, $g_1 \leq h$  for some $h \in \FBLp[X]$.
by \Cref{approximation}(2), there exist $g_1 = G\big(\delta_{x_1}, \ldots, \delta_{x_m}\big) \in \FVL[X]_+$ ($G$ is a lattice expression, $x_1, \ldots, x_m \in X_0$) and $g_2 \in \Hp[X]$ with $\|g_2\| < \vr$, such that $g \leq g_1 + g_2$. %$\|h_1 - h\| < \vr/2$.
%Let $h_2 = g_2 + (h-h_1)$, hence $g \leq h_1 + h_2$, with $\|h_2\| < \vr$. 
If $x_1^*, \ldots, x_n^* \in \cap_{i=1}^m x_i^\perp$, then, for every $j$, $g_1(x_j^*) = 0$, hence $g(x_j^*) \leq g_2(x_j^*)$. Therefore,
$$
\sum_j |g(x_j^*)|^p \leq \|g_2\|^p \|(x_j^*)\|_{p,\textrm{weak}}^p < \vr^p . \qedhere
$$
\end{proof}


\begin{proposition}\label{supremum in FBL}
Suppose $T \in B(Z,X)$ is compact and dual $p$-persistent relative to a dense $X_0 \subset X$. Define $\phi : X^* \to \R : x^* \mapsto \|T^*x^*\|$. Fix a norm one $x_0 \in X_0$ and $\mu > 0$. If $g \in \FBLp[X]$ satisfies $g \geq f = \mu \big| \delta_{x_0} \big| \wedge \phi$, then $g \geq \mu \big| \delta_{x_0} \big|$.
\end{proposition}

\begin{proof}
By scaling, we assume that $c$ in the definition of dual $p$-persistence equals $1$.
Pick a norm one $x^* \in X^*$, and show that $g(x^*) \geq \mu |\langle x^*, x_0\rangle|$. 

Pick $\delta > 0$, and write $g = g_1 + g_2$, where $g_1 = G(\delta_{x_1}, \ldots, \delta_{x_m})$, $G$ being a lattice polynomial, and $\|g_2\| < \delta$.
Find $x_1^*, \ldots, x_n^* \in \cap_{i=0}^m x_i^\perp$ so that $\|(x_j^*)\|_{p,{\textrm{weak}}} \leq 1$, and $\sum_j \alpha_j^p = 1$, where $\alpha_j = \|T^* x_j^*\|$.

Pick $C > \|T\| + \mu + 1$. For each $j$, $g(\alpha_j x^* + C x_j^*) = \alpha_j g_1(x^*) + g_2(\alpha_j x^* + C x_j^*)$, or, in other words,
$$
\big( g(\alpha_j x^* + C x_j^*) \big)_j - \big( \alpha_j \big)_j g_1(x^*) = \big( g_2(\alpha_j x^* + C x_j^*) \big)_j .
$$
The triangle inequality in $\ell_p^n$ gives us:
$$
\Big|\big\| \big( g(\alpha_j x^* + C x_j^*) \big)_j \big\|_p - g_1(x^*) \Big| \leq \big\| \big( g_2(\alpha_j x^* + C x_j^*) \big)_j \big\|_p 
$$
(here we use our scaling: $\big\| \big( \alpha_j \big)_j \big\|_p = 1$). Taking $\|x^*\| = 1$ into account,
$$
\big\| \big( \alpha_j x^* + C x_j^* \big)_j \big\|_{p, {\textrm{weak}}} \leq C \big\| \big( x_j^* \big)_j \big\|_{p, {\textrm{weak}}} + \big\| \big( \alpha_j \big)_j \big\|_p \leq C+1 ,
$$
hence $\big\| \big( g_2(\alpha_j x^* + C x_j^*) \big)_j \big\|_p < \delta (C+1)$, and therefore,
\begin{equation}
\Big|\big\| \big( g(\alpha_j x^* + C x_j^*) \big)_j \big\|_p - g_1(x^*) \Big| < (C+1)\delta .
\label{g g1 close}
\end{equation}

Now note that
$$
f(\alpha_j x^* + C x_j^*) = \mu \alpha_j \big| \langle x^*, x_0 \rangle \big| \wedge \|T^* (\alpha_j x^* + C x_j^*)\| ,
$$
and
$$
\|T^* (\alpha_j x^* + C x_j^*)\| \geq C \|T^* x_j^*\| - \alpha_j \|T^* x^*\| = \alpha_j (C - \|T^*x^*\|) .
$$
By our choice of $C$, $C - \|T^*x^*\| > \mu + 1$, hence $f(\alpha_j x^* + C x_j^*) = \mu \alpha_j \big| \langle x^*, x_0 \rangle \big|$. We clearly have
$$
\big\| \big( g(\alpha_j x^* + C x_j^*) \big)_j \big\|_p \geq \big\| \big( f(\alpha_j x^* + C x_j^*) \big)_j \big\|_p = \mu \big\| \big( \alpha_j \big)_j \big\|_p \big| \langle x^*, x_0 \rangle \big| = \mu \big| \langle x^*, x_0 \rangle \big| ,
$$
so from \eqref{g g1 close}, $g_1(x^*) \geq \mu \big| \langle x^*, x_0 \rangle \big| - (C+1)\delta$. Hence,
$$
g(x^*) \geq g_1(x^*) - \|g_2\| \geq \mu \big| \langle x^*, x_0 \rangle \big| - (C+2)\delta .
$$
To complete the proof, recall that $\delta$ can be arbitrarily small.
\end{proof}


\begin{proof}[Proof of \Cref{failure of fatou}]
We keep the earlier notation: $(Z_k)$ is an NFDR for $T$, $T_k = T|_{Z_k}$, $\phi(x^*) = \|T^* x^*\|$, $\phi_k(x^*) = \|T_k^* x^*\|$. By scaling, we assume that $c$ in the definition of dual $p$-persistence equals $1$.

(1) First we establish the failure of $\lambda$-Fatou for $\FBLp[X]$. Fix $\mu > \lambda \pi_p(T^*)$ and a norm one $x_0 \in X_0$. Let $f_k = \phi_k \wedge \mu \big| \delta_{x_0} \big|$; clearly $f_k \in \FBLp[X]$, and $\|f_k\| \leq \|\phi_k\| \leq \pi_p(T^*)$.
If $g \in \FBLp[X]$, and $g \geq f_k$ for any $k$, then $g \geq \mu \big| \delta_{x_0} \big| \wedge \phi$, hence, by \Cref{supremum in FBL}, $g \geq \mu \big| \delta_{x_0} \big|$, so $\|g\| \geq \mu > \lambda \pi_p(T^*)$.

(2) We show that $\Ip[X]$ is not monotonically bounded by proving that there is no $g \in \Ip[X]$ so that $g \geq \phi_k$ for every $k$.
% The functions $\phi_k$ introduced above form an increasing sequence in $\FBLp[X]_+$. 
Indeed, if such a $g$ exists, then $g \geq \sup_k \phi_k = \phi$. By \Cref{go far to get nothing}, there exist $x_1, \ldots, x_m \in X_0$ so that $\sum_j g(x_j^*)^p < 1/2$ whenever $x_1^*, \ldots, x_n^* \in \cap_i x_i^\perp$ satisfy $\|(x_j^*)\|_{p,{\textrm{weak}}} \leq 1$.
However, by the dual $p$-persistence of $T$, we can find $x_1^*, \ldots, x_n^*$ as above so that $\sum_j \phi(x_j^*)^p > 1$, yielding the desired contradiction.
\end{proof}


\begin{remark}
The proof of \Cref{failure of fatou} shows that if, in addition, $T$ has an SFDR, then $\FBLp[X]$ and $\Ip[X]$ fail the countable $\lambda$-Fatou property and countable monotonic boundedness, respectively.
\end{remark}

In a similar fashion, we show:

\begin{theorem}\label{Ip not Hp general}
Suppose there exists a Banach space $Z$ and a compact, dual $p$-persistent $T \in B(Z, X)$ such that $T^* \in \Pi_p(X^*,Z^*)$. Then $\Hp[X] \supsetneq \Ip[X] \supsetneq \nip[X]$.
\end{theorem}

The proof requires:

\begin{proposition}\label{ideal is closed criterion}
The ideal $\nip[X]$ is closed in $\Hp[X]$ if and only if there exists a constant $C \geq 1$ such that for every $f \in \nip[X]$ there exists $g \in \FBLp[X]$ with $|f| \leq g$ and $\|g\| \leq C \|f\|$.
\end{proposition}

\begin{proof}
Suppose there exists a constant $C$ as above. We show that the sum of any absolutely convergent series in $\nip[X]$ belongs to $\nip[X]$. Indeed, suppose $f = \sum_n f_n$ and $\sum_n \|f_n\| < \infty$ (the sum exists in $\Hp[X]$).
For each $n$, find $g_n \in \FBLp[X]$ such that $|f_n| \leq g_n$ and $\|g_n\| \leq C \|f_n\|$, then $g = \sum_n g_n \in \FBLp[X]$, and $|f| \leq g$, so $f \in \nip[X]$.

Conversely, suppose $C$ as above does not exist. Then, for each $n$, we can find $f_n \in \nip[X]$ so that $\|f_n\| < 2^{-n}$, yet $\|g\| > 2^n$ whenever $g \in \FBLp[X]$ and $|f_n| \leq g$. Then $f = \sum_n f_n \in \Ip[X] \backslash \nip[X]$.
\end{proof}

\begin{proof}[Proof of \Cref{Ip not Hp general}]
Suppose $T \in B(Z,X)$ is compact and dual $p$-persistent relative to a dense $X_0 \subset X$. Define $\phi : X^* \to \R : x^* \mapsto \|T^*x^*\|$. Fix $x_0 \in X_0$.
By \Cref{supremum in FBL}, if $g \in \FBLp[X]$ satisfies $g \geq f = \big| \delta_{x_0} \big| \wedge \phi$, then $g \geq \big| \delta_{x_0} \big|$. So, $\|g\| \geq \|x_0\|$, while $\|f\| \leq \|\phi\| = \pi_p(T^*)$.
In light of \Cref{ideal is closed criterion}, $\Ip[X] \supsetneq \nip[X]$.

Additionally, $\phi \in \Hp[X]$ (due to $T$ being compact with $p$-summing adjoint), while $\phi \notin \Ip[X]$, due to the dual $p$-peristence of $T$ and \Cref{go far to get nothing}. Thus, $\Hp[X] \supsetneq \Ip[X]$.
\end{proof}

\begin{comment}
\section{Some useful classes of operators}

Following up on the previous section, we exhibit classes of ``useful'' operators $T : Z \to X$ which permit us to build examples of failure of monotonic boundedness and of $\lambda$-Fatou property.

{\bf (1)} For $\overline{T} \in B(L_1(0,1), X)$, let $T = \overline{T} \circ id$, where $id : L_{p'} (0,1) \to L_1(0,1)$ is the canonical embedding, and $1/p + 1/p' = 1$.
Throughout, $L_{p'} (0,1)$, with $Z_k = Z_k^{(p')} = \spn[\chi_{A_{ki}} : 1 \leq i \leq 2^k]$, where $A_{ki} = ((i-1) 2^{-k}, i 2^{-k})$, $k \in \N$.
Note that $id^*$ takes $L_\infty(0,1)$ to $L_p(0,1)$ (this is trivial for $p>1$, and is easily checked for $p=1$), so $\cup_k Z_k$ is norming for $\ran(T^*)$, hence $(Z_k)$ is an SFDR. Also, $\pi_p(T^*) \leq \iota_p(T^*) \leq \| \overline{T}\|$.
We should also recall that, by \cite{bourgain-DP}, $T$ is compact if and only if $\overline{T}$ is Dunford-Pettis.

One approach to constructing $\overline{T}$ is to define it on subspaces $Z_k = Z_k^{(1)}$.

\begin{proposition}\label{define operator}
Suppose for each $k$ the operator $S_k : Z_k^{(1)} \to X$ is defined, such that $\sup_k \|S_k\| < \infty$, and $S_m\big|_{Z_k^{(1)}} = S_k$ for $m>k$. Then there exists $S \in B(L_1(0,1),X)$ with $\|S\| = \sup_k \|S_k\|$ and $S\big|_{Z_k^{(1)}} = S_k$ for any $k$.
\end{proposition}

\begin{proof}
For $f \in L_1(0,1)$ define $Sf = \lim_k S P_k f$, where $P_k : L_1(0,1) \to Z_k^{(1)}$ is the conditional expectation. We only need to show the existence of this limit. Note that, for $m > k$,
$S_m P_m f - S_k P_k f = S_m(P_m f - P_k f)$, hence $\|S_m P_m f - S_k P_k f\| \leq \sup_m \|S_m\| \|P_m f - P_k f\|$. Clearly $P_k f \to f$, hence $\lim_{m,k \to \infty} \|P_m f - P_k f\| = 0$.
\end{proof}


{\bf (2)} For the second class, consider finite dimensional spaces $E_k$, and operators $T'_k : E_k \to X$, with $\|T_k'\| \to 0$. Let $Z = (\oplus_{k=1}^\infty E_k)_1$, and $Z_n = (\oplus_{k=1}^n E_k)_1$; define $T : Z \to X : (e_k) \mapsto T_k' e_k$. Then $T$ is clearly compact.
\end{comment}

\section{Failure of Fatou property in function spaces via non-compact operators}\label{no compactness}

In this section, we begin our series of concrete examples of function spaces failing the Fatou property.

Recall that a Banach space $X$ is said to have the Complete Continuity Property (CCP), also referred to as the Compact Range Property, if every operator $T: L_1 \to X$ is Dunford--Pettis (in other words, completely continuous).
% By \cite[Section 7-2]{Talagrand}, the wRNP implies the CCP.

\begin{proposition}\label{no CCP}
    If $X$ fails the CCP, and $1 \leq p < \infty$, then each of the lattices $\FBLp[X]$, $\Ip[X]$, and $\Hp[X]$ fails the $\lambda$-Fatou property, for any $\lambda > 0$.
\end{proposition}

\begin{proof}
   Find $\widetilde{T} \in B(L_1(0,1), X)$ which is not Dunford-Pettis, and let $T = \widetilde{T} \circ id$, where $id : L_{p'} (0,1) \to L_1(0,1)$ is the canonical embedding, and $1/p + 1/p' = 1$. Note that $\pi_p(T^*) \leq \| \widetilde{T}\|$. Moreover, by \cite[Proposition 1]{bourgain-DP}, $T$ fails to be compact. An appeal to \Cref{no compact no fatou} completes the proof.
\end{proof}

\begin{corollary}\label{c0 or L1}
 If $X$ contains either $c_0$ or $L_1(0,1)$, and $1 \leq p < \infty$, then each of the lattices $\FBLp[X]$, $\Ip[X]$, and $\Hp[X]$ fails the $\lambda$-Fatou property, for any $\lambda > 0$.
\end{corollary}
\begin{proof}
    In both cases, $X$ fails the CCP (see, for example, \cite[p. 76]{Talagrand}).
\end{proof}

For $p=1$, \Cref{c0 or L1} was claimed in \cite{ABT}. However, for $X$ non-separable, that proof contains a gap: it relies on $\overline{i}(\FBL[Y])$ being a regular sublttice of $\FBL[X]$ whenever $i : Y \to X$ is an isometric embedding. \Cref{FBLp into Ip} below shows that this is not necessarily correct.


\begin{corollary}\label{duals and Hp}
    If $X$ be a Banach space and $p \in [1,\infty)$, then: 
    \begin{enumerate}
        \item $\Hp[X^*]$ is monotonically bounded if and only if $\Hp[X^*]$ has the $\lambda$-Fatou property for some $($equivalently, any$)$ $\lambda$ if and only if $X^*$ has the wRNP;
        \item $\Ip[X^*]$ has the Fatou property if and only if $X^*$ has the wRNP.
    \end{enumerate}
\end{corollary}

\begin{proof}
    If $X^*$ has the wRNP, then, by \Cref{wRNP_Hp}, $\Hp[X^*]$ is 1-monotonically bounded and therefore has the Fatou property. Since $\Ip[X^*]$ is an ideal in $\Hp[X^*]$, it follows that $\Ip[X^*]$ has the Fatou property as well. 
    
    Conversely, suppose $X^*$ fails the wRNP. Then $X$ contains a copy of $\ell_1$ (see \cite{Musial}), and hence $X^*$ contains a copy of $L_1(0, 1)$ (see \cite{Hag73}).  \Cref{c0 or L1} shows that neither $\Hp[X^*]$ nor $\Ip[X^*]$ has the $\lambda$-Fatou property, for any $\lambda$.
\end{proof}

Specializing to Banach lattices, we obtain:

\begin{corollary}\label{dual of sep lattice}
For a Banach lattice $Y$, the following are equivalent:
\begin{enumerate}
 \item $Y^*$ has the RNP.
 \item $Y^*$ has the wRNP.
 \item For some $p \in [1,\infty)$ and $\lambda \geq 1$, $\FBLp[Y^*]$ has the $\lambda$-Fatou property.
 \item For any $p \in [1,\infty)$ there exists $\lambda \geq 1$ so that $\FBLp[Y^*]$ has the $\lambda$-Fatou property.
 \item For any $p \in [1,\infty)$, $\FBLp[Y^*]$ is 1-monotonically bounded.
\end{enumerate} Furthermore, if $Y$ is separable, then each of the above is equivalent to $Y^*$ being separable.
\end{corollary}
\begin{proof}
 $(1) \Leftrightarrow (2)$ holds for Banach lattices in general, see \cite{GS}. 
 $(5) \Rightarrow (4) \Rightarrow (3)$ is evident, and $(1) \Rightarrow (5)$ follows from \Cref{dual-with-RNP}. It remains to establish $\lnot (2) \Rightarrow \lnot (3)$. As discussed in the proof of the preceding corollary, $\lnot (2)$ implies that $L_1(0, 1)$ embeds into $Y^*$. The desired conclusion therefore follows from \Cref{c0 or L1}.
Finally, for a separable Banach space $Y$, the condition $(1)$ is known to be equivalent to the separability of $Y^*$ (combine Corollary 5.12 and Theorem 5.23 of \cite{BL}). 	
\end{proof}

We conclude this section by considering the case $p = \infty$.
\begin{proposition}
For a Banach space $X$, the following are equivalent:
\begin{enumerate}
    \item $X$ is infinite dimensional.
    \item For any $\lambda > 0$, $\FBL^{(\infty)}[X]$ fails the $\lambda$-Fatou property.
    \item For some $\lambda > 0$, $\FBL^{(\infty)}[X]$ fails the $\lambda$-Fatou property.
\end{enumerate}
\end{proposition}

\begin{proof}
Clearly, $(2) \Rightarrow (3)$. The implication $(3) \Rightarrow (1)$ follows from the fact that, if $\dim X < \infty$, then $\FBL^{(\infty)}[X]$ can be identified with $C(\sphere(X^*))$, which has the Fatou property. For $(1) \Rightarrow (2)$, it suffices to consider the identity operator on $X$ and apply \Cref{no compact no fatou}.
\end{proof}
\begin{comment}
    
\begin{proof}[Proof of \Cref{c0 or L1}]
% Renorming $X$ if necessary, we can assume that $j$ is an iso $L_1$, resp.~$c_0$, embeds into $X$ isometrically. 
In each of the two cases of interest, we shall construct an non-compact operator into $X$, with a $p$-summing inverse. In fact, it suffices to find such an operator with values in $L_1$, resp.~$c_0$.
In the spirit of the preceding section, these will be operators of the form $T = \overline{T} \circ id$, where $\overline{T}$ takes $L_1(0,1)$ to itself or to $c_0$, and $id : L_{p'}(0,1) \to L_1(0,1)$ is the formal identity.

In case of $L_1(0,1)$, we can take $\overline{T}$ to be the formal identity.


 For $c_0$, denote by $(f_i)$ denote its canonical basis. Let $x_{01} = 0$, $x_{11} = f_1$, $x_{12} = -f_1$. Proceed recursively: for $k \in \N$ and $1\leq i \leq 2^k$, let $x_{k+1,2i-1} = x_{ki} + f_{k+1}$, $x_{k+1,2i} = x_{ki} - f_{k+1}$.
 Then, for any $k$ and $i$, $x_{ki} = \sum_{j=1}^k \vr(k,i,j) f_j$, where $\vr(k,i,j) = \pm 1$, and moreover, any $\pm1$ string of length $k$ appears exactly once among strings $(\vr(k,i,j))_{j=1}^k$.

Define $S_k : Z_k^{(1)} \to c_0 : 2^k \chi_{A_{ki}} \mapsto x_{ki}$; clearly this is a contraction, and $S_m\big|_{Z_k^{(1)}} = S_k$ for $m>k$.
By \Cref{define operator}, the sequence $(S_k)$ determines a contraction $\overline{T} : L_1(0,1) \to c_0$. Let $T = \overline{T} \circ id$, then $T^*$ is $p$-summing.

If $T$ were compact, then, by \Cref{weak* continuity}, we would have $\lim_k \|T^* f_k^*\| = 0$, where $(f_k^*)$ is the canonical basis in $\ell_1$. However, note that, for any $i$,
$$
\langle T_k^* f_k^*, \chi_{A_{ki}} \rangle = \langle f_k^*, T_k \chi_{A_{ki}} \rangle = 2^{-k} \langle f_k^*, x_{ki} \rangle = 2^{-k} \vr(k,i,k) .
$$
Consequently,
$$
\|T_k^* f_k^*\| \geq \langle T_k^* f_k^*, \sum_{i=1}^{2^k} \vr(k,i,k) \chi_{A_{ki}} \rangle = 1 ,
$$
hence, for every $k$, $\|T^* f_k^*\| \geq \|T_k^* f_k^*\| \geq 1$.
\end{proof}

\begin{remark}
The above proof shows that $\FBLp[E]$ fails the sequential $\lambda$-Fatou property. Note that a separable lattice has the $\lambda$-Fatou property iff it has the sequential $\lambda$-Fatou property.

\marg{\timur{Can shorten second paragraph, just mention \cite{ABT}}}
Indeed, it is clear that $\lambda$-Fatou implies sequential $\lambda$-Fatou. For the converse, suppose a separable lattice $X$ has the sequential $\lambda$-Fatou property. Suppose $(x_\alpha)_{\alpha \in A}$ is an upward directed net in $X_+$, with $\vee_\alpha \|x_\alpha\| = 1$, and having the supremum $x$.
To show that $\|x\| \leq \lambda$, we use some ideas from \cite{ABT}. Pick a countable dense subset of $(x_\alpha)$, which we shall denote by $(y_i)_{i=1}^\infty$. For $n \in \N$ let $z_n = \vee_{i=1}^n y_i$. It suffices to show that $x = \vee_n z_n$. Indeed, clearly $z_n \leq x$ for any $n$; so, if $x$ is not equal to $\vee_n z_n$, then there exists $z \leq x$ so that $z \neq x$, and $z \geq z_n$ for any $n$.
Note that, for any $\alpha \in A$ and $k \in \N$, there exists $j \in \N$ so that $\|y_j - x_\alpha\| < 2^{-k}$, hence $\| (x_\alpha - z_n)_+ \| < 2^{-k}$ for $n \geq j$. Consequently, $\| (x_\alpha - z)_+ \| < 2^{-k}$. As $k$ is arbitrary, $x_\alpha \leq z$, contradicting $x$ being the least upper bound of $(x_\alpha)$.
\end{remark}
\end{comment}

\section{Failure of Fatou property via dual $p$-persistence: the case of the (pre)dual of the James Tree space}\label{no m-b}\label{dual persistence}\label{james tree}


Consider the dual, and the canonical predual, of the \emph{James Tree space} $JT$. We briefly outline the construction here, and refer the reader to \cite{LiSt} or \cite[Section 15.4]{AK} for more detail.

$JT$ is defined as the completion of all finitely supported double indexed sequences $(a_{ni})_{n \geq 0, 1 \leq i \leq 2^n}$, arranged as a tree in the canonical way -- namely, $(n+1,2i-1)$ and $(n+1,2i)$ are direct descendants of $(n,i)$. The norm of $a = (a_{ni})$ is
\begin{equation}
	\|a\| = \big( \sup \sum_k \big| \sum_{(n,i) \in S_k} a_{ni} \big|^2 \big)^{1/2} ,
	\label{define:JT}
\end{equation}
where the supremum runs over all finite collections of disjoint segments $(S_k)$. % We shall denote the canonical basis of $JT$ by $(e_{ni})$. 
The canonical normalized basis of $JT$, which we denote by $(e_{ni}^*)$, is boundedly complete (see e.g.~\cite[Proposition 15.4.4, Remark 15.4.5]{AK}). Denoting the biorthogonal functionals by $(e_{ni})$, we see that $JT$ is the dual of their span, which we denote by $JT_*$. We refer the reader to \cite{LiSt} for more information on $JT_*$.

\begin{proposition}\label{JT*-Fatou}
For $p \in [2,\infty)$, both $\FBLp[JT_*]$ and $\FBLp[JT^*]$ fail the $\lambda$-Fatou property, for any $\lambda > 0$; also, both $\Ip[JT_*]$ and $\Ip[JT^*]$ fail to be monotonically bounded.
\end{proposition}

\begin{proof}
Let $\mu$ be the canonical probablity measure on the Cantor set $\Delta = \{0,1\}^{\N}$ -- that is, $\mu = \nu^{\otimes \N}$, where $\nu$ is the measure on $\{0,1\}$, assigning $1/2$ to each of the two points.
% We shall build $\overline{T} \in B(L_1(\mu), JT_*)$ (using \Cref{define operator}), and let $T = \overline{T} \circ id$, with $id : L_{p'}(\mu) \to L_1(\mu)$. 
We shall find $\widetilde{T} \in B(L_1(\mu), JT_*)$ and define $T = \widetilde{T} \circ id$, where $id : L_{p'}(\mu) \to L_1(\mu)$. We shall then show that both $T$ and $\kappa T$ are dual $p$-persistent, where $\kappa : JT_* \to JT^*$ is the canonical embedding. Since $T^*$ is clearly $p$-summing, \Cref{failure of fatou} completes the proof.

Consider $R : JT \to C(\Delta) : e_{ni}^* \mapsto h_{ni}$, where $h_{ni}$ is a ``Haar function''.
More specifically, consider the branch leading from the root $(0,1)$ to $(n,i)$, consisting of pairs $(k-1,j_k)$, with $1 \leq k \leq n$ and $j_k = 2j_{k-1} - 1 + a_k$, with $a_k \in \{0,1\}$. For $\sigma = (\sigma_k)_{k \in \N} \in \Delta = \{0,1\}^\N$, let $h_{ni}(\sigma) = 1$ if $\sigma_k = a_k$ for $1 \leq k \leq n$; otherwise, set $h_{ni}(s\sigma) = 0$. This can also be re-stated in terms of the binary expansion of $i-1$.
	
By \cite[p.~92]{LiSt}, $R$ is contractive. Let $\widetilde{T} = R^*\big|_{L_1(\mu)}$, and $x_{ni} = R^* h_{ni}^*$, where $h_{ni}^* = 2^n h_{ni} \in L_1(\mu) \subset C(\Delta)^*$.
%, where $\nu$ is the canonical product measure on $\Delta$ -- that is, for any $\vr_1, \ldots, \vr_n \in \{0,1\}$, the measure of the set of all elements of $\Delta$ having $(\vr_1, \ldots,\vr_n)$ the starting $n$-tuple equals $2^{-n}$.
Then $\langle x_{ni} , e^*_{st} \rangle = 2^n \langle h_{ni} , h_{st} \rangle$ equals to:
	\begin{itemize}
		\item $1$ if $(s,t)$ is a predecessor of $(n,i)$, or $(n,i)$ itself.
		\item $2^{n-s}$ if $(s,t)$ is a successor of $(n,i)$.
		\item $0$ if $(n,i)$ and $(s,t)$ are not on the same branch.
	\end{itemize}
By \cite[Theorem 1]{LiSt} (invoked on p.~92 of that paper, where the notation $y^*_{ni}$ is used), $x_{ni} \in JT_*$; so $\widetilde{T}$ takes values in $JT_*$.
% Clearly $x_{ni} = (x_{n+1,2i-1} + x_{n+1,2i})/2$, giving us (iii). Now we can easily verify (iv).
% :	$$	\sum_{j=1}^{2^k} f_k(x_{kj}^*)^p = 2^{-k} \sum_{j=1}^{2^k} \sum_{i=1}^{2^k} \langle x^*_{kj}, x_{ki} \rangle^p = 1 ,	$$	since $\langle x^*_{kj}, x_{ki} \rangle$ equals $1$ if $i=j$, and $0$ otherwise.
We therefore define $T = \widetilde{T} \circ id$, which coincides with the astriction of $R^* \circ id$ to $JT_*$ ($\kappa T = R^* \circ id$). From the definition of $\widetilde{T}$, 
$$
\langle h_{ni}^* , \widetilde{T}^* e^*_{st} \rangle = \langle h_{ni}^* , h_{st} \rangle \, \, {\textrm{ for \,  any }}  \, \, (n,i) , (s,t) ,
$$
hence $\widetilde{T}^* e^*_{st} = h_{st} \in C(\Delta) \subset L_\infty(\mu)$, and therefore, $T^* e^*_{st} = h_{st} \in L_p(\mu)$. For future use, note that $\|T^* e^*_{st}\| = 2^{-s/p}$.
We shall show the dual $p$-persistence (with constant $1/2$) of $\kappa T : L_1(\mu) \to JT^*$; the statement for $T$ will immediately follow. 

% Let $(x_{ni}) \subset JT_*$ be the tree constructed above, and viewed as a subset of $JT^*$.
Denote by $B$ the set of branches of the James tree (starting with the ``root''). For $b \in B$ and $x^* = (\alpha_{ni}) \in JT$, let $s_b(x^*) = \sum_{(n,i) \in b} \alpha_{ni}$. Note that this sum converges: if not, then the branch $b$ would contain consecutive intervals $I_1, I_2, \ldots$ so that $\inf_k \big| \sum_{(n,i) \in I_k} \alpha_{ni} \big| > 0$, which gives $\|x^*\| = \infty$. In fact, $s_b$ is a norm one linear functional on $JT$.

By \cite{AGM}, $\spn[JT_* , s_b : b \in B]$ is dense in $JT^*$. To establish the dual $p$-persistence (with $X = JT^*$), take $X_0$ to be the linear span of $\{s_b : b \in B\}$, and of finitely supported elements of $JT_*$.
For $F \subset JT^*$, spanned by finitely supported $y_1, ..., y_M \in JT_*$ and $s_{b_1}, ..., s_{b_K}$ ($b_1, \ldots, b_k \in B$), we need to find $x_1^*, \ldots, x_J^* \in F^\perp$ so that $\|(x_j^*)\|_{p,\textrm{weak}} \leq 1$, yet $\sum_j \|T^* x_j^*\|^p \geq 2^{-p}$.
% Now take finitely supported $y_1, ..., y_M \in JT_*$, and $b_1, \ldots, b_k \in B$. We have to find $x_1^*, \ldots, x_J^* \in F^\perp \subset JT$ (where $F = \{y_1, ..., y_M, s_{b_1}, ..., s_{b_K}\}$), for which $\|(x_j^*)\|_{p,\textrm{weak}} \leq 1$, and $\sum_j \|T^* x_j^*\|^p \geq 2^{-p}$.
To this end, find $s > 2 \log K$ so that $\langle e_{st}^*, y_i \rangle = 0$ for every $i \in \{1, \ldots, M\}$ and $t \in \{1, \ldots, 2^s\}$.
Let $x_1^*, \ldots, x_J^*$ be an enumeration of those $e_{st}^*$ which lie outside of the branches $b_1, \ldots, b_K$, then $J \geq 2^{s-1}$.
As no two nodes $(s,t)$ (for fixed $s$) can belong to the same segment, we have
$\| \sum_t \alpha_t e^*_{st} \big\| = \big( \sum_t |\alpha_t|^2 \big)^{1/2}$.
Consequently, if $\sum_j |\beta_j|^{p'} \leq 1$, then
$$
\| \sum_j \beta_j x_j^* \| = \big( \sum_j |\beta_j|^2 \big)^{1/2} \leq 1 
$$
(here we rely on $p' \leq 2$), hence $\| (x_j^* )\|_{p, {\textrm{weak}}} = 1$. On the other hand, $\|T^* x_j^*\| = 2^{-s/p}$, hence $\sum_j \|T^* x_j^*\|^p = J 2^{-s} \geq 1/2$, since $J \geq 2^{s-1}$.
\end{proof}

\begin{corollary}
 If $p \in [2, \infty)$ and $n \geq 1$ is odd, then $\FBLp[JT^{(n)}]$ fails the $\lambda$-Fatou property, for any $\lambda > 0$.
\end{corollary}

\begin{proof}
 Write $n = 2m-1$, and proceed by induction on $m$. \Cref{JT*-Fatou} gives us the base case of $m = 1$. For the inductive step, recall that, being a dual space, $JT^{(2m - 1)}$ is contractively complemented in $JT^{(2m + 1)}$. 
 Therefore, by \cite{OiTTT} (see also \cite{OiTTT-cor}), $\FBLp[JT^{(2m - 1)}]$ embeds lattice isometrically into $\FBLp[JT^{(2m + 1)}]$, as a regular sublattice. Clearly, the $\lambda$-Fatou property passes down to closed regular sublattices.
\end{proof}

\begin{remark}
    Since $JT^*$ has the wRNP (because $JT$ does not contain a copy of $\ell_1$), $\Ip[JT^*]$ has the Fatou property by \Cref{duals and Hp}. However, $\Ip[JT^*]$ is not monotonically bounded, by \Cref{JT*-Fatou}.
\end{remark}

\begin{remark}\label{bourgain}
\begin{comment}
It is well-known that uniformly bounded $X$-valued martingales correspond to bounded linear operators $T: L_1(0,1) \to X$. By \cite[Theorem 1]{bourgain-DP}, $T: L_1(0,1) \to X$ is Dunford--Pettis if and only if the corresponding martingale is Cauchy in the Pettis norm.
%$$
%\| \xi \|=\sup \left \{\int |\langle \xi(t), x^*\rangle | \, dt: \; x^* \in \ball(X^*) \right\}.
%$$
Moreover, by \cite[Proposition 8]{bourgain-DP}, if a martingale $(\xi_n)$ is Cauchy in the Pettis norm, then for every weak$^*$-null sequence $(x_k^*)$, we have
$$
\lim_{k \to \infty} \sup_n \int|\langle \xi_n(t),x_k^*\rangle| \, dt = 0.
$$
Since $(\xi_n)$ is uniformly bounded, the above is equivalent to
\begin{equation}\label{non-bonsai}
  \lim_{k \to \infty} \sup_n \Big(\int |\langle \xi_n(t), x_k^*\rangle|^p \, dt \Big)^{1/p}=0, \quad 1 < p < \infty.  
\end{equation}
\end{comment}
It is known (see, for example, \cite[Chapter 7]{Talagrand}) that the wRNP implies the CCP, so the operator $\widetilde{T}: L_1 \to JT_*$ constructed above is Dunford-Pettis. Since $id: L_{p'} \to L_1$ is weakly compact, $T = \widetilde{T} \circ id : L_{p'} \to JT_*$ is compact. For this reason, the results of \Cref{operators:general} tell us nothing about $\Hp[JT_*]$.
% we cannot conclude anything concerning $\Hp[JT_*]$.
% Since $JT^*$ has the wRNP, every uniformly bounded martingale with values in $JT^*$, and hence also with values in $JT_*$, satisfies \eqref{non-bonsai}. Therefore, one cannot construct a bonsai bush in either $JT_*$ or $JT^*$, and must resort to espalier trees.
\end{remark}

\begin{comment}
    

\begin{corollary}
 If $n$ is odd, then $\FBLp[JT^{(n)}]$ fails the $\lambda$-Fatou property, for any $\lambda > 1$.
\end{corollary}

\begin{proof}
It is known (see e.g. \cite{AGM} or \cite{LiSt}) that $JT^*/JT_* = \ell_2(\Gamma)$, where $\Gamma$ is the set of all infinite branches (having the cardinality of the continuum). Consequently, $JT^*$ is isomorphic to $JT^* \oplus \ell_2(\Gamma)$.
Moreover, $JT^{**}$ is isomorphic to $JT \oplus \ell_2(\Gamma)$. From this, we conclude that $JT^{(n)} = JT^*$ whenever $m$ is odd (and also $JT^{(n} = JT^*$ whenever $n$ is even, but this is not relevant here). Therefore, it suffices to establish the failure of Fatou for $\FBLp[JT_*]$ and $\FBLp[JT^*]$.

In light of \Cref{FBLp no Fatou}, it suffices to build an espalier tree in $JT^*$. To this end, we ``recycle'' $x_{ni} \in JT_*$ constructed earlier; we view $JT_*$ ($JT$) as canonically embedded into $JT^*$ (resp.~$JT^{**}$). It remains to find $x_{kj}^* \in JT$ ($k \in \N, 1 \leq j \leq 2^k$), so that:
\begin{enumerate}[(i)]
\item For any $k$, $\| \big((x_{kj}^*)_{j=1}^{2^k}\big) \|_{2, {\textrm{weak}}} = 1$.
\item The sequence of $2^k$-tuples $\big((x_{kj}^*)_{j=1}^{2^k}\big)_k$ is $2$-weak$^*$ null in $JT^{**}$ -- that is, for any $x^{**} \in JT^*$, $\lim_k \big( \sum_i | \langle x^{**} , x^*_{ki} \rangle |^p \big)^{1/p} = 0$. 
% \item For any $n$ and $i$, $x_{n+1,2i-1} + x_{n+1,2i} = 2x_{ni}$.
\item For $n \in \N$, let $f_n = \big( 2^{-n} \sum_{i=1}^{2^n} \big| \delta_{x_{ni}} \big|^p \big)^{1/p}$; then $\sum_{j=1}^{2^n} f_n(x_{nj}^*)^p \geq 1$.
\end{enumerate}
Indeed, by the $2$-convexity of $\FBL^{(2)}[JT^*]$, we have $\|f_n\| \leq 1$. As in the proof of \Cref{bonsai tree}, we show that the sequence $(f_n)$ is increasing. The proof is completed by invoking \Cref{limit}.

For $k \in \N$ and $1 \leq j \leq N_k := 2^k$, consider $x_{kj}^* = 2^{-k/2} \sum_{i=(j-1)2^k+1}^{j2^k} e_{2k i}^* \in JT$. In other words, $x_{kj}^*$ is built as follows: it has zero entries on all levels except $2k$. On the level $2k$, it has $2^{-k/2}$ on all branches passing through the $j$-th node on the $k$-th level, zero otherwise.

To establish (i), suppose  $\sum_j |\alpha_j|^2 \leq 1$. In \eqref{define:JT}, each of the disjoint segments $S_k$ involved meets the support of $\sum_j \alpha_j x_{kj}^*$ at most once. Consequently, $\|\sum_j \alpha_j x_{kj}^*\| = \big(2^k \sum_j (2^{-k/2} |\alpha_j|)^2 \big)^{1/2} \leq 1$.
	
Now we show (ii). By \cite[Proposition 3.8]{AGM}, $JT^*$ is the closed linear span of elements canonical unit vectors $e_{ni} \in JT_*$ and $\sigma \in \Gamma$, where $\sigma$ denotes the summation along a branch. In light of (i), it suffices to establish the limit from (ii) for $x^{**}$ of this type only.
This, however, is straightforward.

To verify (iii), note that 
	$$
	\sum_{j=1}^{2^k} f_k(x_{kj}^*)^p = 2^{-k} \sum_{j=1}^{2^k} \sum_{i=1}^{2^k} \langle x^*_{kj}, x_{ki} \rangle = 1 ,
	$$
	since $\langle x^*_{kj}, x_{ki} \rangle$ equals $1$ if $i=j$, and $0$ otherwise.
\end{proof}

\end{comment}
\begin{comment}
\section{The case of $\FBL^{(\infty)}[E]$}\label{p=infty}

The case of $p=\infty$ is different from that of $p$ finite.

\begin{proposition}
For a Banach space $E$, the following are equivalent:
\begin{enumerate}
    \item $E$ is infinite dimensional.
    \item For any $\lambda > 0$, $\FBL^{(\infty)}[E]$ fails the $\lambda$-Fatou property.
    \item For some $\lambda > 0$, $\FBL^{(\infty)}[E]$ fails the $\lambda$-Fatou property.
\end{enumerate}
\end{proposition}

\begin{proof}
Clearly, $(2) \Rightarrow (3)$. The implication $(3) \Rightarrow (1)$ follows from the fact that, if $\dim E < \infty$, then $\FBL^{(\infty)}[E]$ can be identified with $C(\sphere(E^*))$, which has the Fatou property. It remains to prove $(1) \Rightarrow (2)$.

 Fix an infinite dimensional $E$ and $\lambda > 1$. Fix $x_0 \in \sphere(E)$ (the unit sphere of $E$) for which there exists $y_0 \in \sphere(E^*)$ s.t. $\langle y_0,x_0 \rangle = 1$ (this is possible by Bishop-Phelps-Bolobas Theorem).
 Find a separable $Y \subset X^*$, containing $y_0$, and a separable $X_0 \subset X$ which norms $Y$, and contains $x_0$. Let $(x_i)_{i=1}^\infty$ be a dense subset of $\sphere(X_0)$. For $n \in \N$ let $f_n = \lambda |\delta_{x_0}| \wedge (\vee_{i=1}^n |\delta_{x_i}|)$.
 Clearly $(f_n)$ is an increasing sequence in $\ball(\FBL^{(\infty)}[E])_+$, bounded above by $\lambda |\delta_{x_0}|$. We shall show that, if $f$ is an upper bound for $(f_n)$, then $\|f\| \geq \lambda$.
 
 Suppose, for the sake of contradiction, that an upper bound $f$ as above satisfies $\|f\| < \lambda$. Pick $\vr \in (0, (\lambda - \|f\|)/2)$; there exist $z_1, \ldots, z_N \in E$, and a non-negative lattice expression $F$ so that, for any $x^* \in E^*$, (i) $F((\langle x^*,z_j\rangle)_{j=1}^N) \leq (\|f\| + \vr) \|x^*\|$, and (ii) for any $n$, $f_n(x^*) \leq F((\langle x^*,z_j\rangle)_{j=1}^N) + \vr \|x^*\|$.
 Condition (ii) implies that, for $x^* \in Y$, we have $\lambda |\langle x^*,x_0 \rangle| \wedge \|x^*\| \leq F((\langle x^*,z_j\rangle)_{j=1}^N) + \vr \|x^*\|$.
 
 For any $y \in Y \cap x_0^\perp \cap( \cap_j  z_j^\perp )$, we have
 $$
 F((\langle y_0 + y,z_j\rangle)_{j=1}^N) = F((\langle y_0,z_j\rangle)_{j=1}^N) \leq \|f\|+\vr ,
 $$
 hence
 $$
 \vee_n f_n(y_0 + y) = \lambda \wedge \|y_0+y\| \leq F((\langle y_0 + y,z_j\rangle)_{j=1}^N) + \vr \|y_0+y\| \leq \|f\| + \vr + \vr \|y_0+y\| .
 $$
 Selecting $y$ so that $\|y_0+y\| = \lambda$, we conclude that $\lambda \leq \|f\| + \vr(\lambda + 1)$. As $\vr$ can be arbitrarily small, conclude that $\lambda \leq \|f\|$.
\end{proof}
\end{comment}


\section{Embedding of $\nip[X]$ into $\Hp[X]$}\label{FBLp vs Ip vs Hp}\label{Ip into Hp}
In \Cref{la-tr}, we showed that if $X$ is a dual Banach space with the RNP, and $1 \leq p < \infty$, then $\Hp[X] = \nip[X]$. This statement fails for general $X$.


\begin{theorem}\label{Ip not HP}
In the following cases, $\Hp[X] \supsetneq \Ip[X] \supsetneq \nip[X]$:
\begin{enumerate}
    \item $X$ contains $c_0$, $1 \leq p < \infty$.
    \item $X = JT_*$ or $X = JT^*$, $2 \leq p < \infty$.
    \item $X$ is a subspace of $L_1(0,1)$ which fails the RNP while possessing the Bounded Approximation Property $($BAP$)$, $p=1$.
\end{enumerate}
\end{theorem}

Note that the spaces $X$ appearing in the three cases above are distinct. For instance, $JT^*$ has the wRNP, since its predual doesn't contain a copy of $\ell_1$. By \cite[Theorem 7.4.8]{Bour}, $c_0$ cannot be embedded into $JT^*$, \emph{a fortiori} into $JT_*$.

By \Cref{Ip not Hp general}, it suffices to find a compact dual $p$-persistent $T \in B(Z,X)$ with $T^* \in \Pi_p(X^*,Z^*)$. Such a $T$ is constructed on a case by case basis.

\begin{proof}[Proof of \Cref{Ip not HP}(2)]
The operator $T$ constructed in the proof of \Cref{JT*-Fatou} is dual $p$-persistent, and has $p$-summing adjoint. Moreover, it is compact, in light of \Cref{bourgain}.
\end{proof}


\subsection{Proof of \Cref{Ip not HP}(1)}
Build $T : \ell_1 \to X$ with the desired properties.

 It is easy to see (by tracking extreme points) that there exists a quotient map $q_k : \ell_1^{2^k} \to \ell_\infty^k$. Then $q_k^* : \ell_1^k \to \ell_\infty^{2^k}$ is an isometric embedding.
 Let ${\boldsymbol{\kappa}}_k = \big( \pi_p (q_k^*)\big)^{-1} = \big( \pi_p (I_{\ell_1^k})\big)^{-1}$. Let $S_k = {\boldsymbol{\kappa}}_k q_k$, then $\pi_p(S_k^*) = 1$.
 Represent $c_0 = (\oplus_k \ell_\infty^k)_0$ and $\ell_1 = (\oplus_k \ell_1^{2^k})_1$, and let $S = \oplus_k S_k : \ell_1 \to c_0$.

 \begin{lemma}\label{kappa-k}
 For $p \in [1,\infty)$, there exists $\gamma = \gamma_p$, so that the inequalities $\gamma^{-1} \leq \sqrt{k} {\boldsymbol{\kappa}}_k \leq \gamma$ hold for any $k \in \N$.
 \end{lemma}

\begin{proof}
In \cite{Gar}, $p$-summing norms of diagonal operators are computed (up to a constant). For the reader's convenience, here we outline a proof of $\pi_p (I_{\ell_1^k}) \sim \sqrt{k}$.

For the lower estimate recall that $\ell_1^k$ contains a subspace $G$ with $\dim G \geq ck$ ($c$ is a universal constant), and $d_{BM}(G, \ell_2^k) \leq 2$ (here and below, $d_{BM}(\cdot , \cdot )$ stands for the Banach-Mazur distance). Then
$$
\pi_p (I_{\ell_1^k}) \geq \pi_p(I_G) \geq \frac12 \pi_p(I_{\ell_2^k}) \geq c_p \sqrt{k} ,
$$
where the last inequality utilizes \cite[Corollary 4.13]{DJT} ($c_p$ depends on $p$ only).

On the other hand, the cotype $2$ constants of the $\ell_1^k$ spaces are uniformly bounded, hence the $1$- and $2$-summing norms are equivalent for operators on $\ell_1^k$ \cite[Corollary 11.16]{DJT}. Now recall that $\pi_2(I_F) = \sqrt{\dim F}$ for any $F$.
\end{proof}
 
 \begin{lemma}\label{p-summing norm}
% In the above notation, $\pi_p(S^*) = 1$.
For Banach spaces $X_1, Y_1, X_2, Y_2, \ldots$, define $X = (\oplus_k X_k)_1, Y = (\oplus_k Y_k)_\infty$. 
Suppose $u_k \in B(X_k,Y_k)$ are given, and let $u = \oplus_k u_k \in B(X,Y)$. Then, for $p \in [1,\infty]$, $\pi_p(u) = \sup_k \pi_p(u_k)$.
 \end{lemma}

\begin{proof}
The equality for $p=\infty$, and the inequality $\pi_p(u) \geq \sup_k \pi_p(u_k)$, are easy to verify. So it suffices to show that $\pi_p(u) \leq 1$ whenever $\sup_k \pi_p(u_k) = 1$ for finite $p$.
To this end, pick $x_1, \ldots, x_N \in X$ with $\|(x_j)\|_{p,{\textrm{weak}}} < 1$. We need to establish that $\big( \sum_j \|u x_j\|^p \big)^{1/p} \leq 1$. 
Write $x_j = (x_{jk})$, where $x_{jk} \in X_k$. We can and do assume that the vectors $x_j$ have finite support, that is, $x_{jk} = 0$ for $k > K$. Then
\begin{equation}
\sum_k \big\| \sum_j \gamma_j x_{jk} \big\| \leq 1 \, \, {\textrm{  whenever  }} \, \, \|(\gamma_j)_j\|_{\ell_{p'}} \leq 1 .
\label{eq:bound on sum}
\end{equation}
Represent $\{1, \ldots, N\}$ as a disjoint union of sets $S_k$ ($1 \leq k \leq K$) so that, for $j \in S_k$, $\|u x_j\| = \|u_k x_{jk}\|$. Then
$$
\big( \sum_j \|u x_j\|^p \big)^{1/p} = \Big( \sum_k \sum_{j \in S_k} \|u_k x_{jk}\|^p \Big)^{1/p} .
$$
Find nonnegative $(\alpha_k)$ so that $ \|(\alpha_k)_k\|_{\ell_{p'}} = 1$ and 
$$
\Big( \sum_k \sum_{j \in S_k} \|u_k x_{jk}\|^p \Big)^{1/p} =
\sum_k \alpha_k \Big( \sum_{j \in S_k} \|u_k x_{jk}\|^p \Big)^{1/p} .
$$
The inequality $\pi_p(u_k) \leq 1$ implies
$$
\big( \sum_j \|u x_j\|^p \big)^{1/p} \leq \sum_k \alpha_k \| (x_{jk})_{j \in S_k} \|_{p, \textrm{weak}} .
$$
It remains to show that the right-hand side does not exceed $1$.

For each $k$, find $\beta_j$ ($j \in S_k$) so that $\|(\beta_j)_{j \in S_k}\|_{\ell_{p'}} = 1$ and 
$$
\| \sum_{j \in S_k} \beta_j x_{jk} \| = \| (x_{jk})_{j \in S_k} \|_{p, \textrm{weak}} .
$$
Then
$$
\sum_k \alpha_k \| (x_{jk})_{j \in S_k} \|_{p, \textrm{weak}} = \sum_k \alpha_k \| \sum_{j \in S_k} \beta_j x_{jk} \| = \sum_k \| \sum_{j \in S_k} \gamma_j x_{jk} \| ,
$$
where, for $j \in S_k$, $\gamma_j = \alpha_k \beta_j$. Then 
$\|(\gamma_j)_j\|_{\ell_{p'}} = 1$, since $\gamma_j = \alpha_k \beta_j$ for $j \in S_k$, $\|(\alpha_k)_k\|_{\ell_{p'}} = 1$, and $\|(\beta_j)_{j \in S_k}\|_{\ell_{p'}} = 1$ for every $k$. The result now follows from \eqref{eq:bound on sum}.
\begin{comment}
Clearly $\pi_p(S^*) \geq \sup_k \pi_p(S_k^*) = 1$. To show the converse, pick $x_1, \ldots, x_N \in \ell_1$ with $\|(x_j)\|_{p,{\textrm{weak}}} < 1$. We need to establish that $\big( \sum_j \|S^* x_j\|^p \big)^{1/p} \leq 1$. 
To this end, write $x_j = (x_{jk})$, where $x_{jk}$ belongs to the $k$-th summand of $\ell_1$ (a copy of $\ell_1^k$). Then
\begin{equation}
\sum_k \big\| \sum_j \gamma_j x_{jk} \big\| \leq 1 \, \, {\textrm{  whenever  }} \, \, \sum_j |\gamma_j|^{p'} \leq 1 .
\label{eq:bound on sum}
\end{equation}
We can and do assume that the vectors $x_j$ have finite support, that is, $x_{jk} = 0$ for $k > K$. Represent $\{1, \ldots, N\}$ as disjoint union of sets $S_k$ ($1 \leq k \leq K$) so that, for $j \in S_k$, $\|S^* x_j\| = \|S^*_k x_{jk}\|$. Then
$$
\big( \sum_j \|S^* x_j\|^p \big)^{1/p} = \Big( \sum_k \sum_{j \in S_k} \|S^*_k x_{jk}\|^p \Big)^{1/p} .
$$
Find $(\alpha_k)$ so that $\sum_k |\alpha_k|^{p'} = 1$ and 
$$
\Big( \sum_k \sum_{j \in S_k} \|S^*_k x_{jk}\|^p \Big)^{1/p} =
\sum_k \alpha_k \Big( \sum_{j \in S_k} \|S^*_k x_{jk}\|^p \Big)^{1/p} .
$$
As noted above, $\pi_p(S_k^*) = 1$, hence
$$
\big( \sum_j \|S^* x_j\|^p \big)^{1/p} \leq \sum_k \alpha_k \| (x_{jk})_{j \in S_k} \|_{p, \textrm{weak}} .
$$
It remains to show that the right side doesn't exceed $1$.

For each $k$, find $\beta_j$ ($j \in S_k$) so that $\sum_{j \in k} |\beta_j|^{p'} = 1$ and 
$$
\| \sum_{j \in S_k} \beta_j x_{jk} \| = \| (x_{jk})_{j \in S_k} \|_{p, \textrm{weak}} .
$$
Then
$$
\sum_k \alpha_k \| (x_{jk})_j \|_{p, \textrm{weak}} = \sum_k \alpha_k \| \sum_{j \in S_k} \beta_j x_{jk} \| = \sum_k \| \sum_{j \in S_k} \gamma_j x_{jk} \| ,
$$
where, for $j \in S_k$, $\gamma_j = \alpha_k \beta_j$. Then
$$
\sum_j |\gamma_j|^{p'} = \sum_k \alpha_k^{p'} \sum_{j \in S_k} |\beta_j|^{p'} = 1 ,
$$
and the result follows from \eqref{eq:bound on sum}.
\end{comment}
\end{proof}

\begin{comment}
\begin{remark}
The proof of \Cref{p-summing norm} actually yields the following: suppose, for $k \in \N$, $U_k \in B(\ell_1^{N_k},Z_k)$. % , with $\sup_k \pi_p(U_k^*) < \infty$. 
Let $U = \oplus_k U_k$ act from $\ell_1 = (\oplus_k \ell_1^{N_k})_1$ to $(\oplus_k Z_k)_0$. Then $\pi_p(U^*) = \sup_k \pi_p(U_k^*)$.
\end{remark}
\end{comment}

Now let $T = jS$, where $j : c_0 \to X$ is an embedding (we make our life more convenient by renorming $X$ to make $j$ isometric).
We have already established that $T$ is compact; by \Cref{p-summing norm}, $T^*$ is $p$-summing. The dual $p$-persistence of $T$ (with $X_0 = X$) follows from:

\begin{lemma}\label{p-sum lower est}
 There exists a positive ${\boldsymbol{\gamma}}$ $($which may depend on $p)$ with the following property: if $j : c_0 \to X$ is an isometric embedding, and $E$ is a subspace of $X^*$ of finite codimension, then $\pi_p(S^* j^*|_E) \geq {\boldsymbol{\gamma}}$.
\end{lemma}



\begin{proof}
The quotient map $j^* : X^* \to \ell_1$ admits a lifting $R : \ell_1 \to X^*$, of norm less than $2$. Let $F = R(\ell_1)$, then $\|x^*\|/2 \leq \|j^* x^*\| \leq \|x^*\|$ for every $x^* \in F$. 

Recall that we write $S^* = \oplus_k S_k^*$, where $S_k^*$ is defined on $G_k = \ell_1^k$, the ``$k$-th summand'' of $\ell_1$. Then $F_k = R(G_k)$ is $2$-isomorphic to $\ell_1^k$. Clearly $\dim E \cap F_k \geq k - \codim E$.
By \cite[Theorem 2.1]{GTT}, for $k$ large enough $F_k \cap E$ contains a subspace $H$, $c$-isomorphic to $\ell_1^m$, with $m \geq k/c$ ($c > 1$ is a constant). Hence,
$$
\pi_p \big(S^* j^*|_E\big) \geq \pi_p \big( S^* j^* |_H \big) \geq \frac1{2c} {\boldsymbol{\kappa}}_k \pi_p \big(I_{\ell_1^m}\big) = \frac1{2c} \frac{{\boldsymbol{\kappa}}_k}{{\boldsymbol{\kappa}}_m} \geq {\boldsymbol{\gamma}} ; 
$$
here, we use the estimate ${\boldsymbol{\kappa}}_k^{-1} \sim \sqrt{k}$ from \Cref{kappa-k}.
\end{proof}


\subsection{Proof of \Cref{Ip not HP}(3)}

Our first lemma is known, we sketch a proof for convenience.

\begin{lemma}\label{compute nu1}
For any $u \in B(\ell_\infty^N,Z)$, we have $\pi_1(u) = \iota_1(u) = \nu_1(u) = \sum_j \|u \delta_j\|$, where $(\delta_j)_{j=1}^N$ is the canonical basis of $\ell_\infty^N$.
\end{lemma}

\begin{proof}
Clearly, $\|(\delta_j)_j\|_{1,{\textrm{weak}}} = 1$, so $\nu_1(u) \geq \iota_1(u) \geq \pi_1(u) \geq \sum_j \|u \delta_j\|$.
On the other hand, denote by $(\sigma_j)$ the canonical basis of $\ell_1^N$. 
Let $D = {\textrm{diag}}(\|u \delta_j\|) : \ell_\infty^N \to \ell_1^N$ (that is, $D \delta_j = \|u \delta_j\| \sigma_j$) and $v : \ell_1^N \to Z : \sigma_j \mapsto u \delta_j/\|u \delta_j\|$.
Then $v$ is contractive, and $u = vD$, hence $\nu_1(u) \leq \|v\| \|D\| = \sum_j \|u \delta_j\|$.
\end{proof}

By the definition of the BAP, we can find finite rank maps $R_n \in B(X)$ so that $R_m R_n = R_n$ when $m >n$ (in other words, $R_m\big|_{\ran R_n} = I_{\ran R_n}$), $R_n \to I_X$ point-norm, and $\sup_n \|R_n\| = \lambda < \infty$.
We shall construct an operator $\widetilde{T} \in B(L_1(0,1),X)$ such that, for the formal identity $id : L_\infty(0,1) \to L_1(0,1)$, the operator $T = \widetilde{T} \circ id$ is compact and dual $1$-persistent relative to $\cup_n \ran(R_n)$. Clearly, both $T$ and $T^*$ will be $1$-summing.

By \cite{bourgain-DP}, there exists $\widetilde{T} : L_1(0,1) \to X$ which is Dunford-Pettis, but not representable. As noted in that paper, $T$ will be compact. It remains to show the dual $1$-persistence.
Throughout, we use the notation $R_n^\circ = I - R_n$.

\begin{lemma}\label{1-summing}
There exists $c_0 > 0$ such that, for any $n$, $\pi_1(R_n^\circ T) \geq c_0$.
\end{lemma}

\begin{proof}
Suppose otherwise. %Note that $\overline{T}$, being non-representable, cannot have finite rank, so $R_n^\circ T \neq 0$ for any $n$. 
Then there exist $n_1 < n_2 < \ldots$ such that $\pi_1(R_{n_k}^\circ T) < 2^{-k}/\lambda$. 
As the domain of $R_{n_k}^\circ T$ is $L_\infty(0,1)$, the $1$-summing and $1$-integral norms of this operator coincide. Hence, by \cite[Proposition 5.27]{DJT},
$$
\nu_1(R_{n_{k+1}} R_{n_k}^\circ T) \leq \|R_{n_{k+1}}\| \iota_1(R_{n_k}^\circ T) < 2^{-k} .
$$
%Recall that the operator $R_{n_1} T$, being of finite rank, is also nuclear.
Note that, for any $K$, $R_{n_{K+1}} - R_{n_K} = R_{n_{K+1}} R_{n_K}^\circ$. By induction, 
$$
R_{n_1} + R_{n_2} R_{n_1}^\circ + \ldots + R_{n_{K+1}} R_{n_K}^\circ = R_{n_{K+1}} .
$$
%As $R_m^\circ\big|_{\ran(R_n)} = 0$ for $m > n$, we conclude that
%$$
%Tx = R_{n_1} Tx + R_{n_2} R_{n_1}^\circ Tx + \ldots % + R_{n_{K+1}} R_{n_K}^\circ T + R_{n_{K+1}}^\circ R_{n_K}^\circ T . 
%$$
%for any $x$ from $\cup_n \ran(R_n)$, which is dense in $X$.
Hence, the sum $R_{n_1} T + R_{n_2} R_{n_1}^\circ T + \ldots$ is Cauchy in $\nu_1( \cdot )$. Since the ideal of nuclear operators is complete, this sum must converge, necessarily to $T$, as $R_{n_k}T \to T$ in the operator norm. Here we use the fact $T$ is compact and $R_{n_k} \to I_X$ uniformly on compact subsets of $X$.
Thus, $T$ is nuclear, hence $\widetilde{T}$ is representable \cite[C3]{DeFl}. This is the desired contradiction.
\end{proof}

\begin{proof}[Proof of \Cref{Ip not HP}(3), conclusion]
We need to show that, for every $n$, we can find $x_1^*, \ldots, x_N^* \in \ran(R_{n-1})^\perp$ such that $\|(x_j^*)\|_{1,{\textrm{weak}}} \leq 1$, while $\sum_j \|T^* x_j^*\| > c$, where $c$ is a universal constant.

By \Cref{1-summing}, $\pi_1(R_n^\circ T) \geq c_0$ for some $c_0 > 0$. Recall that $R_n^\circ T$ takes values in $X \subset L_1(0,1)$. A discretization argument shows that there exists a finite rank contraction $U : L_1(0,1) \to \ell_1^N$ so that $\pi_1(S T) > c_0/2$, where $S = U R_n^\circ$. We can identify $S$ with a tuple $(x_1^*, \ldots, x_N^*)$ via $Sx = ( \langle x_j^* , x \rangle )_j$.  Note that $S$ vanishes on $\ran(R_{n-1})$, so $x_j^*$ vanish on $\ran(R_{n-1})$, and $\|(x_j^*)\|_{1,{\textrm{weak}}} = \|S\| \leq \lambda + 1$.
As $ST$ takes $L_\infty(0,1)$ to $\ell_1^N$, we have $\pi_1(ST) = \iota_1(ST) = \iota_1(T^*S^*)$ (see \cite[Chapter 5]{DJT}).
%By \cite[5.9]{DeFl} (see also \cite[Exercise 5.10]{DeFl}), $\nu_1(ST) = \nu_1(T^*S^*)$. 
% As $L_1(0,1)$ is contractively complemented in its bidual, we can view $T^* S^*$ as acting from $l_\infty^N$ to $L_1(0,1)$. Then $\nu_1(T^*S^*) = \iota_1(T^*S^*) = \pi_1(T^*S^*)$.
Denote by $(\delta_j)$ the canonical basis of $\ell_\infty^N$, then $S^* \delta_j = x_j^*$ hence $T^* S^* \delta_j = T^* x_j^*$. By \Cref{compute nu1}, $\sum_j \|T^* x_j^*\| = \iota_1(T^* S^*)$, which is known to exceed $c_0/2$.
\end{proof}


\section{The closure of $I_{w^*}$ in $\| \cdot \|_\infty$} \label{sup-norm}
The situation changes drastically if we consider the $\sup$ norm. For the sake of brevity, we denote $\ball(X^*)$, equipped with the weak$^*$ topology, by $K$. 
All the function spaces we are dealing with -- in particular, $I_{w^*}[X] \subset \nip[X] \subset \Ip[X] \subset \Hp[X]$ are linear subspaces of $C(K)$; we equip the latter with its usual $\sup$ norm $\| \cdot \|_\infty$. % By positive homogeneity, we can extend these functions unambiguously to the whole of $X^*$. Below, the notation $\ball(\Hp[X])$ etc.~ refers to the unit balls in the $\| \cdot \|_p$ norm.

\begin{proposition}\label{denseness}
For any Banach space $X$, and $1 \leq p \leq \infty$, % $\ball(I_{w^*}[X])$ is dense in $\ball(\Hp[X])$.
$$ \overline{\ball(I_{w^*}[X])}^{\| \cdot \|_\infty} = \overline{\ball(\Ip[X])}^{\| \cdot \|_\infty} = \ball(\Hp[X]) . $$
Here, $\ball(I_{w^*}[X])$ consists of those $f \in I_{w^*}[X]$ for which $\|f\|_p \leq 1$, where $\| \cdot \|_p$ is the norm inherited from $\Hp[X]$.
\end{proposition}


\begin{proof}
For $p = \infty$, this result is contained in \cite{OiTTT}. Henceforth, we assume $p \in [1,\infty)$. %In keeping with the usual notation, we define $\delta_x : K \to \R : x^* \mapsto \langle x^*,x \rangle$ ($x \in X$).

Clearly $\ball(\Ip[X]) \subset \ball(\Hp[X])$, and, since the $\| \cdot \|_\infty$ limits do not increase the $\| \cdot \|_p$ norm, $\overline{\ball(\Ip[X])}^{\| \cdot \|_\infty} \subset \ball(\Hp[X])$. Our goal is to establish $ \overline{\ball(I_{w^*}[X])}^{\| \cdot \|_\infty} \supset \ball(\Hp[X])$.

Denote by $A$ the set of all positively homogeneous functions $g \in C(K)$, for which there exist $x_1, \ldots, x_N \in X$ so that $\sum_{i=1}^N \|x_i\|^p \leq 1$, and $|g| \leq \big( \sum_i |\delta_{x_i}|^p \big)^{1/p}$. 
Further, let $A^p = \{h \in C(K) : |h|^{1/p} \in A\}$. We view $A^p$ as a subset of $W$ -- the sublattice of $C(K)$ consisting of functions $g$ which satisfy $g(sx^*) = s^p g(x^*)$, for any $s \geq 0$ and $x^* \in X^*$. %Again, $Z_p$ is a sublattice of $C(K)$.
In other words, $h \in W$ lies in $A^p$ iff there exist $x_1, \ldots, x_N \in X$ so that $\sum_{i=1}^N \|x_i\|^p < 1$, and $|h| \leq \sum_i  |\delta_{x_i}|^p$. %By \cite{OiGEO}, $A$ is a dense subset of $\ball(\Ip[X])$.

First deal with positive functions: show that, if $f \in \Hp[X]_+$, and $f^p$ doesn't belong to the $\| \cdot \|_\infty$-closure $A^p$, then $\|f\|_p \geq 1$. 

It is easy to see that $A^p$ is solid and convex. Therefore, by \cite{OiTu}, there exists $\phi \in W^*_+$ such that $\langle \phi , f^p \rangle \geq 1$, and % $\langle \phi , g \rangle \leq 1$ for any $g \in A^p$.
\begin{equation}
\langle \phi , g \rangle \leq 1 \, \, {\textrm{for  any  }} \, \, g \in A^p .
\label{eq:action on Ap}
\end{equation}
Denote by $j$ the embedding $W \to C(K)$. Then $j^* : M(K) \to W^*$ is a weak$^*$ continuous quotient map ($M(K) = C(K)^*$ is the space of finite Radon signed measures on $K$), hence we can find $\mu \in M(K)$ so that $j^* \mu = \phi$, and $\|\mu\| = \|\phi\|$.
Write $\mu = \mu_+ - \mu_-$, then $j^* \mu_+ \geq \phi$, hence $\|\mu_+\| \geq \|\phi\| = \|\mu\|$. On the other hand, $M(K)$ is an AL-space, hence $\|\mu\| = \|\mu_+\| + \|\mu_-\|$. This is only possible if $\mu_- = 0$, so $\phi$ is implemented by a positive Radon measure $\mu$.

Fix $\vr > 0$. For each $x^* \in K$, there exists an open set $U(x^*) \subset K$ so that $|f(x^*)^p - f(y^*)^p| \leq \vr/(2\|\mu\|)$ whenever $y^* \in U(x^*)$. Moreover, for this $U(x^*)$ there exist $x_1, \ldots, x_M \in X$ (depending on $x^*$) so that
$$
\big\{ y^* \in K : \max_i | \langle y^*, x_i \rangle | < 1 \big\} \subset U(x^*) .
$$
The set on the left is convex, and therefore, we can assume that $U(x^*)$ is convex for any $x^*$. Such sets form an open cover of $K$, from which we can extract a finite subcover, say $U(x_1^*), \ldots, U(x_L^*)$.
Then $U(x_1^*)$, $U(x_2^*) \backslash U(x_1^*)$, $\cdots$, $U(x_L^*) \backslash \cup_{j=1}^{L-1} U(x_j^*)$ are disjoint Borel sets covering $K$. By tossing out sets of measure $0$, we obtain disjoint Borel sets $V_1, \ldots, V_N \subset K$ s.t. (i) $\mu(V_i) > 0$ for any $i$, and $\mu(K \backslash (\cup_i V_i)) = 0$, and (ii) for any $i$, $|f(y^*)^p - f(z^*)^p| < \vr/\|\mu\|$ whenever $y^*, z^* \in \overline{{\textrm{conv}}(V_i)}$. 

Let $t_i = \mu(V_i)$ and $\nu_i = t_i^{-1} \chi_{V_i} \mu$ (so $\nu_i$ is a probability measure supported on $V_i$). By \cite[Theorem 1.1]{Phelps}, any $V_i$ possesses a barycenter -- that is, a (unique) point $x_i^* \in \overline{{\textrm{conv}}(V_i)}$ so that $\langle x_i^*, x \rangle = \int \langle x^*, x \rangle d\nu_i(x^*)$ for any $x \in X$. 

Plugging $g = |\delta_x|^p$ ($x \in \sphere(X)$) into \eqref{eq:action on Ap}, we obtain:
\begin{align*}
1 
& 
\geq \big\langle \mu, |\delta_x|^p \big\rangle = \int_K | \langle x^*, x \rangle|^p \, d\mu(x^*) = \sum_i t_i \int | \langle x^*, x \rangle|^p \, d\nu_i(x^*)
\\ &
\geq \sum_i t_i \Big| \int \langle x^*, x \rangle  \, d\nu_i(x^*) \Big|^p = \sum_i t_i \big| \langle x_i^*, x \rangle \big|^p  = \sum_i \big| \langle t_i^{1/p} x_i^*, x \rangle \big|^p 
\end{align*}
or equivalently, $\|(t_i^{1/p} x_i^*)\|_{p,{\textrm{weak}}} \leq 1$.

On the other hand, for each $i$ and $x^* \in V_i$, we have $f(x^*)^p \leq f(x_i^*)^p + \vr/\|\mu\|$, hence
\begin{align*}
1 
&
\leq \int f(x^*)^p \, d\mu(x^*) = \sum_i t_i \int_{V_i} f(x^*)^p \, d\nu_i(x^*) 
\\ &
\leq \sum_i t_i \Big( f(x_i^*)^p + \frac{\vr}{\|\mu\|} \Big) = \sum_i f\big( t_i^{1/p} x_i^* \big)^p + \vr ,
\end{align*}
and so, $\|f\|_p^p \geq \sum_i f\big( t_i^{1/p} x_i^* \big)^p \geq 1-\vr$, and we are done, since $\vr>0$ is arbitrary.

For a general $f \in \ball(\Hp[X])$, write $f = f_+ - f_-$. We have shown that, for any $\vr > 0$, there exists a positive $g \in A \subset \ball(I_{w^*}[X])$ so that $\| |f| - g \|_\infty < \vr$.
Now let $h = f_+ \wedge g - f_- \wedge g$. Clearly $|h| = f_+ \wedge g + f_- \wedge g \leq g$, hence $h \in \ball(I_{w^*}[X])$. Also, $$|f - h| = |f_+ - f_+ \wedge g| + |f_- - f_- \wedge g| \leq ||f|-g|,$$ so $\|f - h\|_\infty < \vr$.
\end{proof}


\section{Embedding of $\FBLp[X]$ into $\nip[Y]$ and $\FBLp[Y]$}\label{FBLp into Ip}

In \cite{LaTr}, it was shown that, if $X$ is a Banach space with an infinite-dimensional separable quotient, then, for $p \in [1, \infty)$, $\FBLp[X]$ is a proper subspace of $\Ip[X]$. We obtain the following strengthening of this result. 

\begin{theorem}\label{very irregular}
Let $X$ and $Y$ be infinite-dimensional Banach spaces, and let $i: X \to Y$ be an isometric embedding. Then, for $p \in [1,\infty)$, % there exist an increasing sequence $(f_n) \subset \FBLp[X]_+$, $f \in \FBLp[X]_+$ and $g \in \nip[X]_+$ such that $f = \vee_n f_n$ in $\FBLp[X]_+$, while $\overline{i} f \neq \overline{i} g = \vee_n \overline{i} f_n$ in $\nip[Y]$. Consequently, 
$\overline{i}(\FBLp[X])$ is not a regular sublattice of $\nip[Y]$.
\end{theorem}

In the particular case of $X=Y$, we have

\begin{corollary}
$\FBLp[X]$ coincides with $\Ip[X]$ if and only if either $X$ is finite dimensional, or $p=\infty$.
\end{corollary}

\begin{proof}[Proof of \Cref{very irregular}]
For $k \in \N$, find $E_k \subset X^*$ which is $2$-isomorphic to $\ell_2^k$.
Further, find a contractive embedding $w_k : E_k \to \ell_\infty^{N_k}$ with $\|w_k^{-1}\| < 2$. Extend $w_k$ to a contraction $v_k : X^* \to \ell_\infty^{N_k}$.
By \cite{OjP}, there exists a contraction $u_k : \ell_1^{N_k} \to X$ so that, for every $x^* \in E_k$, $\|u_k^* x^*\| \geq \|x^*\|/2$.

%Find a decreasing sequence $(\alpha_k) \subset (0,1)$ so that $\alpha_k \sqrt{\log k} \nearrow \infty$. 
Define $T : \ell_1 = (\oplus_k \ell_1^{N_k})_1 \to X : (\xi_k) \mapsto \sum_k {\boldsymbol{\kappa}}_k u_k \xi_k$, where ${\boldsymbol{\kappa}}_k = \big (\pi_p(I_{\ell_2^k})\big )^{-1}$. By \cite[Corollary 4.13]{DJT}, ${\boldsymbol{\kappa}}_k \sim k^{-1/2}$, so $T$ is a compact contraction.
We shall show that $T$ is dual $p$-persistent with respect to $X$. Consider a finite dimensional $F \subset X$, and find $k \geq 2 \dim F$. Then $\dim (E_k \cap F^\perp) \geq k/2$, and 
\begin{align*}
\pi_p(T^*|_{F^\perp}) \geq \kappa_k \pi_p(u_k^*|_{E_k \cap F^\perp}) 
&
\geq \frac{\kappa_k}2 \pi_p(I_{E_k \cap F^\perp}) \geq \frac{\kappa_k}4 \pi_p\big(I_{\ell_2^{\dim (E_k \cap F^\perp)}}\big) 
\\ &
\geq \frac{\kappa_k}{4 \kappa_{\lceil k/2 \rceil}} \geq c , 
\end{align*}
with $c$ being a universal constant.

Now let $\phi(x^*) = \|T^* x^*\| \in H_{w^*}[X]$, let $T_n$ be the restriction of $T$ to $(\oplus_{k=1}^n \ell_1^{N_k})_1$, and denote $\phi_n(x^*) = \|T_n^* x^*\|$ (which clearly belongs to $\FBLp[X]$, and even to $\FVL[X]$). Pick $x \in X$ with $\|x\| > 1$; note that $\big| \delta_x \big|$ is not majorized by $\phi$ since $\phi(x^*) \leq 1$ for every norm-one $x^* \in X^*$. 
Let $f_n = \phi_n \wedge \big| \delta_x \big|$, in light of \Cref{supremum in FBL}, $\vee_n^{\FBLp[X]} f_n = \big| \delta_x \big|$. 

On the other hand, $\overline{i} \big| \delta_x \big| = \big| \delta_{ix} \big|$, and $\overline{i} \phi_n = \phi_n' \in \FBLp[Y]$, with $\phi_n'(y^*) = \|T_n^* i^* y^*\|$. Also, $\overline{i} \phi = \phi'$, with $\phi'(y^*) = \|T^* i^* y^*\|$. Thus $\phi' = \vee_n \phi_n'$ pointwise, so 
$$
\vee_n^{\Ip[Y]} \phi_n' = \phi' \wedge \big|\delta_{ix}\big| = \overline{i} \big( \phi \wedge \big|\delta_x\big| \big) \neq \overline{i} \big|\delta_x\big| , % \qedhere
$$
due to the injectivity of $\overline{i}$.
\end{proof}

% \section{Embedding of $\FBLp[X]$ into $\FBLp[Y]$}

It was asked in \cite{AD} (and, indirectly, in \cite{OiTTT} and \cite{OiTTT-cor}) whether, for an isometric embedding $i : X \to Y$, $\overline{i}(\FBLp[X])$ is necessarily a regular sublattice of $\FBLp[Y]$. The answer is known to be positive when $p=\infty$, and when $i(X)$ is complemented in $Y$. Below we show that the answer is negative in general.

\begin{theorem}\label{complem l2}
Suppose $i : X \to Y$ is an isometric embedding, and there exists $c > 0$ so that, for every $s$, there exists $F_s \subset X^*$, $c$-isomorphic to $\ell_2^s$ and $c$-complemented in $X^*$, via a weak$^*$ continuous projection $Q_s$.
Suppose, furthermore, that $\lim_s s^{-1/2} \pi_p(Q_s i^*) = 0$. Then $\overline{i}(\FBLp[X])$ is not a regular sublattice of $\FBLp[Y]$.
\end{theorem}

\begin{remark}\label{weak-star complementability}
Standard local reflexivity results \cite{OjP} show that the following are equivalent:
\begin{itemize}
    \item $X^*$ containing copies of $\ell_2^n$ uniformly complementably.
    \item $X^*$ containing copies of $\ell_2^n$, uniformly complemented by weak$^*$-continuous projections (in fact, an existing projection can be ``turned into'' a weak$^*$ continuous one).
    \item $X$ containing copies of $\ell_2^n$ uniformly complementably.
\end{itemize}
\end{remark}

\begin{proof}
Find $s_1 < s_2 < \ldots$ such that $s_k^{-1/2} \pi_p(Q_{s_k} i^*) < 2^{-k}$. For convenience, we use notation $E_k = F_{s_k}$ $P_k = Q_{s_k}$. Also, let $R_k \in B(X)$ be such that $R_k^* = P_k$.

As in the proof of \Cref{very irregular}, we can find contractions $u_k \in B(\ell_1^{N_k}, X)$ so that $\|u_k^* x^*\| \geq \|x^*\|/2$ for every $x^* \in E_k$. Consider the operator
$$
T : \ell_1 = \big( \oplus_k \ell_1^{N_k} \big)_1 \to X : (\xi_k) \mapsto \sum_k \gamma_k R_k u_k \xi_k ,
\, \, {\textrm{ where }} \, \gamma_k = \big( c \pi_p(I_{\ell_2^{s_k}}) \big)^{-1}.
$$
Note that $\gamma_k^{-1} \sim \sqrt{s_k}$, and hence $T$ is a compact contraction. Just as in the proof of \Cref{very irregular}, we show that $T$ is dual $p$-persistent relative to $X$.

Let $\phi(x^*) = \|T^* x^*\| \in H_{w^*}[X]$, let $T_n$ be the restriction of $T$ to $(\oplus_{k=1}^n \ell_1^{N_k})_1$, and denote $\phi_n(x^*) = \|T_n^* x^*\| \in \FBLp[X]$. Choose $x_0 \in X$ with $\|x_0\| > 1$, then $\big| \delta_{x_0} \big|$ is not majorized by $\phi$. Define $f_n = \big| \delta_{x_0} \big| \wedge \phi_n \in \FBLp[X]$ and $f = \big| \delta_{x_0} \big| \wedge \phi$.

By \Cref{supremum in FBL}, $\vee_n^{\FBLp[X]} f_n = \big| \delta_{x_0} \big|$. We shall show that $\vee_n^{\FBLp[Y]} \overline{i} f_n = \overline{i} f = \big| \delta_{i x_0} \big| \wedge \overline{i} \phi$; this would complete the proof, as $\overline{i}$ is injective.

Clearly, it suffices to show that $\overline{i} \phi \in \FBLp[Y]$. For any $n \in \N$ and $y^* \in Y^*$, $|\phi_{n}(i^* y^*) - \phi_{n-1}(i^*y^*)| \leq \gamma_n \|u_n^* P_n i^* y^*\|$, hence
$$
\big\| \overline{i} (\phi_{n} - \phi_{n-1}) \big\| \leq \gamma_n \|u_n\| \pi_p(P_n i^*) \leq \frac{\pi_p(P_n i^*)}{\pi_p(I_{E_n})} \leq c 2^{-n} .
$$
Consequently, $(\overline{i} \phi_n)$ is a Cauchy sequence, which converges to $\overline{i} \phi$ pointwise, hence also in $\FBLp[Y]$.
\end{proof}

Applying the preceding proposition, we obtain:

\begin{corollary}\label{scipt L-infty l1}
Suppose $X$ contains $\ell_2^n$ uniformly complementably, and $i : X \to Y$ is an isometric embedding. Then $\overline{i}(\FBLp[X])$ is not a regular sublattice of $\FBLp[Y]$ if one of the following holds:
\begin{enumerate}
    \item $Y$ is a ${\mathcal L}_\infty$ space, and $1 \leq p < \infty$. 
    \item $Y$ is a ${\mathcal L}_1$ space, and $2 \leq p < \infty$. 
\end{enumerate}
\end{corollary}

This corollary is applicable, for instance, when $X$ contains an infinite dimensional K-convex complemented subspace \cite[Theorem 19.3]{DJT}.

\begin{proof}
Keep the notation of \Cref{complem l2}: subspaces $\ell_2^s \sim F_s \subset X^*$ are complemented in $X^*$ via weak$^*$ continuous projections $Q_s$, with $\sup_s \|Q_s\| < \infty$.
If $Y$ is a $\mathcal L_\infty$ space, then $Y^*$ is a $\mathcal L_1$ space \cite{LiRo}, hence, by Grothendieck Theorem, $\pi_1(Q_s i^*) \sim \|Q_s i^*\|$, which implies $\sup_s \pi_1(Q_s i^*) < \infty$.
The case of $Y$ being a $\mathcal L_1$ space is similar, and relies on $\sup_s \pi_2(Q_s i^*) < \infty$.
\end{proof}

We also obtain:

\begin{theorem}\label{l q uniformly}
Suppose $p \in [2,\infty)$, $q \in \{\infty\} \cup [1,2)$ and $Y$ contains copies of $\ell_q^n$ uniformly. Then $Y$ contains a subspace $X$ so that $\overline{i}(\FBLp[X])$ is not a regular sublattice of $\FBLp[Y]$, where $i : X \to Y$ stands for the canonical embedding. 
\end{theorem}

The ``Maurey-Pisier + Krivine'' Theorem (see e.g. \cite[Theorem 6]{Maurey} or \cite[Theorem 13.2]{MS}) tells us that $Y$ contains $\ell_q^n$ uniformly if $q$ is the supremum of all numbers $r$ for which $Y$ has type $r$; likewise, $Y$ contains $\ell_\infty^n$ uniformly iff it doesn't have non-trivial cotype.

For a proof, we need \Cref{complem l2}, as well as a series of auxiliary results.

\begin{proposition}\label{liftings}
There is a function $\psi : (0,1/2) \to (0,1)$ % with $\lim_{t \to 0} \psi(t) = ???$
with the following property.
If $p \in [2,\infty)$, $u \in B(W,\ell_2^n)$ is a contraction with $\pi_p(u) \geq c \sqrt{n}$, then, for $n$ large enough, there exists $E \subset \ell_2^n$ with $\dim E \geq \psi(c) n$, and $v \in B(E,W)$ with $\|v\| \leq 1/\psi(c)$, so that $uv = I_E$.
\end{proposition}

Note that, in the above situation, $\pi_p(u) \leq \pi_2(u) \leq \|u\| \pi_2(I_{\ell_2^n}) = \sqrt{n}$.

\begin{proof}
It clearly suffices to supply a proof for $p=2$.

Take $u$ as above. By trace duality, there exists $T \in B(\ell_2^n,W)$ with $\pi_2(T) \leq c^{-1} \sqrt{n}$ so that $\nu_1(uT) = {\textrm{tr}}(uT) = n$. In fact, we can assume that $uT$ is diagonal in an orthonormal basis $(\delta_i)_{i=1}^n$, with entries $\lambda_1 \geq \ldots \geq \lambda_n \geq 0$.
We have $\sum_i \lambda_i = n$, and $\sum_i \lambda_i^2 = \pi_2(uT)^2 \leq \pi_2(T)^2 \leq c^{-2} n$. We claim that the set $S = \{i : c^4 \leq \lambda_i \leq c^{-4}\}$ has cardinality of at least $c^4 n/2$. 

Indeed, let $S_+ = \{i : \lambda_i > c^{-4}\}$, then $|S_+| c^{-8} < \sum_{i \in S_+} \lambda_i^2 < c^{-2} n$, so $|S_+| < c^6 n$. Further, 
$$
\sum_{i \in S_+} \lambda_i \leq \big( \sum_{i \in S_+} \lambda_i^2 \big)^{1/2} |S_+|^{1/2} \leq c^{-1} \sqrt{n} \cdot c^3 \sqrt{n} = c^2 n .
$$
Finally, $\sum_{i : \lambda_i < c^4} \lambda_i \leq c^4 n$, so $\sum_{i \in S} \lambda_i \geq (1 - c^2 - c^4) n > n/2$, yielding $|S| \geq c^4 n/2$.

Now let $E_0 = \spn[\delta_i : i \in S]$. For any $x \in E_0$, $\|Tx\| \geq \|uTx\| \geq c^4 \|x\|$. Our goal is to find $E \subset E_0$, of dimension proportional to $n$, so that $\|T|_E\| \leq c^{-6}$.
To do this, find $e_1 \in E_0$ with $\|e_1\| = 1$, $\|T e_1\| = \|T|_{E_0}\|$. Let $E_1 = E_0 \cap e_1^\perp$, and find $e_2 \in E_1$ with $\|e_2\| = 1$, $\|T e_2\| = \|T|_{E_1}\|$, then let $E_2 = E_0 \cap e_1^\perp \cap e_2^{\perp}$. 
After performing $k$ steps ($k = \lceil c^9 n \rceil$), we obtain orthonormal vectors $e_1, \ldots e_k$. Clearly, $\|(e_i)_{i=1}^k\|_{2,{\textrm{weak}}} \leq 1$, hence $k \|Te_k\|^2 \leq \sum_{i=1}^k \|T e_i\|^2 \leq \pi_2(T)^2 \leq c^{-2} n$. From this, $\|Te_k\|^2 \leq c^{-2} n/\lceil c^9 n \rceil$.

Let $E = E_0 \cap (\cap_{i=1}^{k-1} e_i^\perp)$. Then $\dim E \geq \dim E_0 - k + 1 \geq \big( c^4/2 - c^9) n$, and $\|T|_E\| \leq \|Te_k\| \leq c^{-6}$.
\end{proof}


It is known (see e.g.~\cite{FLM} or \cite[Section 5]{MS}) that there exists $c > 0$ so that, for $1 \leq r < 2$, $\ell_r^{m(n)}$ contains a $2$-isomorphic copy of $\ell_2^n$, with $m(n) \leq cn$. Consequently, there is a quotient map $q_n : \ell_{r'}^{m(b)} \to F_n$, where $F_n$ is $2$-isomorphic to $\ell_2^n$.

\begin{corollary}\label{bad liftings}
In the above situation, $\limsup_n n^{-1/2} \pi_2(q_n) = 0$.
\end{corollary}

\begin{proof}
Suppose otherwise; then, by \Cref{liftings}, there exist $n_1 < n_2 < \ldots$, for which we can find $E_k \subset F_{n_k}$ and liftings $v_k : E_k \to \ell_{r'}^{m(n_k)}$ with $q_{n_k} v_k = I_{E_k}$, $\|v_k\| \leq C$, and $\dim E_k \geq n_k/C$, where $C$ is an absolute constant.
Consequently, $\ell_{r'}^{m(n_k)}$ contains a $2C$-isomorphic copy of a Hilbert space, of dimension proportional to $m(n_k)$. This, however, is impossible, see \cite[Theorem 1.1]{BDGJN} or \cite[Example 3.1]{FLM}.
\end{proof}

The following two results are part of Banach space lore.

\begin{lemma}\label{adjoint of embedding}
Suppose $j : G \to Z$ is an isometric embedding, and there exists a projection $P \in B(Z)$, with $\ran (P) = j(G)$.
Then, for any $z^* \in \ran(P^*)$, $\|j^* z^*\| \geq \|z^*\|/\|P\|$.
\end{lemma}

\begin{proof}
For $z^*$ as above, 
$$
\|z^*\| = \|P^* z^*\| = \sup_{z \in \ball(Z)} \langle P^* z^* , z \rangle = \sup_{z \in \ball(Z)} \langle z^* , Pz \rangle .
$$
For $z \in \ball(Z)$, $Pz = jg$ for some $g \in G$ with $\|g\| \leq \|P\|$. Thus,
$$
\|z^*\| \leq \sup_{g \in G, \|g\| \leq \|P\|} \langle z^* , jg \rangle = \sup_{g \in G, \|g\| \leq \|P\|} \langle j^*z^* , g \rangle = \|P\| \|j^*z^*\| .
\qedhere
$$
\end{proof}

\begin{lemma}\label{complementation}
Suppose $W$ is a finite dimensional subspace of an infinite dimensional Banach space $Z$. Then for every $\vr > 0$ there exists a finite codimensional subspace $Z_0 \subset Z$, containing $W$, and a projection from $Z_0$ onto $W$, of norm less that $1 + \vr$.
\end{lemma}

\begin{proof}[Sketch of a proof]
Find a contraction $T : W \to \ell_\infty^N$, so that $\|Tw\| \geq \|w\|/(1+\vr)$ for any $w \in W$. Extend $T$ to a contraction $S : Z \to \ell_\infty^N$. We claim that $Z_0 = W + \ker S$ has the desired properties.
Indeed, any $z \in Z_0$ has a unique expansion $z = w + x$, with $w \in W, x \in \ker S$. Then
$$
\|w+x\| \geq \|S(w+x)\| = \|Sw\| \geq \frac{\|w\|}{1+\vr} .
$$
Thus, the map $w+x \mapsto w$ determines a projection from $Z_0$ onto $W$, of norm not exceeding $1+\vr$.
\end{proof}

\begin{lemma}\label{fin codim subspaces}
Suppose $Z$ contains $\ell_q^n$'s uniformly with distortion $C$. Then every finite-codimensional subspace of $Z$ contains $\ell_q^n$'s uniformly with the same distortion.
\end{lemma}
\begin{proof}[Sketch of a proof]
Let $W$ be a subspace of $Z$ with $\dim Z/W = m$. Fix $n \in \N$, and show that $W$ contains an $n$-dimensional subspace $C$-isomorphic to $\ell_q^n$. To this end, find $E \subset Z$ such that $\dim E = n(m+1)$, and there exists a contraction $T : \ell_q^{n(m+1)} \to E$ with $\|T^{-1}\| \leq C$.
Denote the canonical basis of $\ell_q^{n(m+1)}$ by $(\delta_j)_{j=1}^{n(m+1)}$. For $1 \leq k \leq n$, the $T(\spn[\delta_{(k-1)(m+1) + j} : 1 \leq j \leq m+1])$ has dimension $m+1$, hence it meets $W$.
Consequently, for each such $k$ there exists a norm one $e_k \in \spn[\delta_{(k-1)(m+1) + j} : 1 \leq j \leq m+1]$ with $Te_k \in W$. Clearly $\spn[Te_k : 1 \leq k \leq n]$ is $C$-isomorphic to $\ell_q^n$.
%
% Choose a subspace $E \subset Z$ that is $C$-isomorphic to $\ell_q^{n(m+1)}$, and let $T: E \to \ell_q^{n(m+1)}$ be an isomorphism. Note that $\dim E/(W\cap E) \leq m$, hence $T(W \cap E)$ contains an isometric copy of $\ell_q^n$ (to see this, split the basis into $n$ blocks of size $m+1$ and choose a norm-one vector from the span of each block intersected with $T(W \cap E)$). Hence, $W$ contains $\ell_q^n$'s uniformly with the same distortion.
\begin{comment}
Suppose $Z$ contains $\ell_q^n$'s with distortion $C$, and $W$ is a subspace of $Z$ with $\dim Z/W = m$. We shall show that, for every $D>C$, $W$ contains $\ell_q^n$'s with distortion $D$. To this end, fix $\vr \in (0,1/(3Cn))$ so that
\begin{equation}
\label{choice of vr}
\frac{C(1+3n\vr)}{1-3Cn\vr} < D .
\end{equation}
A standard ``volume comparison'' argument (see e.g.~\cite{AuSz}) shows that $\ball(Z/W)$ contains $\xi_1, \ldots, \xi_M$ ($M \leq (3/\vr)^m$) so that, for any $\xi \in \ball(Z/W)$, there exists $j$ with $\|\xi - \xi_i\| \leq \vr$.

Now pick $N \geq 2n+M$, and find a contraction $u : \ell_q^N \to Z$ with so that $\|u^{-1}\| \leq C$. Let $(z_j)_{j=1}^N$ be the images of the canonical basis of $\ell_q^N$ under $u$. Then, for any $\alpha_1, \ldots, \alpha_N \in \R$,
\begin{equation}
\big( \sum_j |\alpha_j|^q \big)^{1/q} \geq \|\sum_j \alpha_j z_j \| \geq C^{-1} \big( \sum_j |\alpha_j|^q \big)^{1/q} .
\label{equivalence}
\end{equation}

Represent $\{1, \ldots, N\}$ as a disjoint union $\cup_{i=1}^M S_i$ so that, if $j \in S_i$, then $\|Q z_j - \xi_i\| \leq \vr$ (here $Q : Z \to Z/W$ is the quotient map).
The set $S_i$ contains $\lfloor |S_i|/2 \rfloor$ pairs of disjoint indices $j$. Thus, there exist at least
$$
\sum_{i=1}^M \Big\lfloor \frac{|S_i|}2 \Big\rfloor \geq \sum_{i=1}^M \frac{|S_i|-1}2 = \frac12 \sum_{i=1}^M |S_i| - \frac{M}2 = \frac{N-M}2 \geq n
$$
disjoint pairs of numbers from $\{1, \ldots, N\}$, so that both elements of each pair lie in the same set $S_i$. Denote such pairs as $(j(k,1), j(k,2))$, with $1 \leq k \leq n$.

For $1 \leq k \leq n$, let $x_k = 2^{-1/q} \big( z_{j(k,1)} - z_{j(k,2)} \big)$. Then $\|Q x_k\| \leq 2\vr$, hence there exists $w_k \in W$ so that $\|x_k - w_k\| < 3 \vr$. Then, for any $\alpha_1, \ldots, \alpha_n \in \R$, 
\begin{equation}
\label{elements of W}
\sum_k \alpha_k w_k = 2^{-1/q} \sum_k \alpha_k x_k + \sum_k \alpha_k (w_k - x_k)
\end{equation}
Note that, if $\sum_k |\alpha_k|^q = 1$, then
$$
\big\| \sum_k \alpha_k (w_k - x_k) \big\| \leq \sum_k |\alpha_k| \cdot \max_k \|w_k - x_k\| \leq 3 n \vr .
$$
From \eqref{equivalence}, $1 \geq \|\sum_k \alpha_k x_k\| \geq C^{-1}$.
Applying the triangle inequality to \eqref{elements of W}, we get:
$$
\big\| \sum_k \alpha_k w_k \big\| \leq n \vr + 2^{-1/q} \big\| \sum_k \alpha_k x_k \big\| \leq 1 + 3n\vr ,
$$
and also,
$$
\big\| \sum_k \alpha_k w_k \big\| \geq 2^{-1/q} \big\| \sum_k \alpha_k x_k \big\| - 3n\vr \geq \frac1C - 3n\vr ,
$$
We have thus found that $W$ contains $\ell_q^n$ with distortion 
$$
\frac{1+3n\vr}{C^{-1}-3n\vr} ,
$$
which is less than $D$, by \eqref{choice of vr}.
\end{comment}
\end{proof}

\begin{proof}[Proof of \Cref{l q uniformly}]
Find increasing sequences $(m(s)), (n(s))$ so that, for every $s$, there is an embedding $j_s : G_s \to \ell_q^{m(s)}$, where $G_s$ is $2$-isomorphic to $\ell_2^{n(s)}$, and $\pi_2(j_s^*) \leq 2^{-s} \sqrt{n(s)}$.
For $q \in (1,2)$, this is possible by \Cref{bad liftings}, while for $q \in \{1,\infty\}$, we only need to invoke Grothendieck Theorem.

Now suppose $Y$ contains $\ell_q^n$ uniformly. By \Cref{fin codim subspaces}, there exists $C$ so that any finite codimensional subspace of $Y$ contains $\ell_q^n$ with distortion $C$. 
We construct $X$ recursively. First find $F_1 \subset Y$, $C$-isomorphic to $\ell_q^{m(1)}$. Denote the embedding of $F_1$ into $Y$ by $a_1$. 
Find $E_1 \subset F_1$ which is $2C$-isomorphic to $\ell_2^{n(1)}$, and such that $\pi_2(b_1^*) \leq C 2^{-1} \sqrt{n(1)}$, where $b_1$ is the embedding of $E_1$ into $F_1$.
Find a finite codimensional $Y_1 \subset Y$, containing $E_1$, and such that there is a projection $P_1$ from $Y_1$ onto $E_1$, of norm less that $1 + 2^{-1}$.

Find $F_2 \subset \ker P_1 = Y_1'$, $C$-isomorphic to $\ell_q^{m(2)}$; denote its embedding into $Y$ by $a_2$.
Find $E_2 \subset F_2$ which is $2C$-isomorphic to $\ell_2^{n(2)}$, and such that $\pi_2(b_2^*) \leq C 2^{-2} \sqrt{n(2)}$, where $b_2$ is the embedding of $E_2$ into $F_2$.
Find a finite codimensional $Y_2 \subset Y_1$, containing $E_1 + E_2$, and such that there is a projection $P_2$ from $Y_2$ onto $E_1 + E_2$, of norm less that $1 + 2^{-2}$.

Proceed in the same manner. Then $X = \spn[E_k : k \in \N]$ contains $E_k$ uniformly complementably: $P_k$ implements a projection from $X$ onto $E_1 + \ldots + E_k$, of norm not exceeding $1 + 2^{-k}$, giving rise to the natural projections $Q_k = (I - P_{k-1}) P_k$ from $X$ onto $E_k$. 
Then $\sup_k \|Q_k\| \leq 4$, and, for every $k$, $H_k = Q_k^*(X^*)$ is $8C$-isomorphic to $\ell_2^{n(k)}$.
On the other hand, $iQ_k j_k = a_kb_k$, where $j_k : E_k \to X$ is the canonical embedding. Therefore, $\pi_2(j_k^* Q_k^* i^*) \leq \pi_2(b_k^*) \leq C 2^{-k} \sqrt{n(k)}$.
However, as $\|Q_k\| \leq 4$, $\|j_k^* x^*\| \geq \|x^*\|/4$ for any $x^* \in H_k$, by \Cref{adjoint of embedding}. Consequently, $\pi_2(Q_k^* i^*) \leq 4 \pi_2(j_k^* Q_k^* i^*) \leq C 2^{2-k} \sqrt{n(k)}$. It remains to recall \Cref{complem l2}.
\end{proof}



% \section{Possible directions for further research}\label{s:further research}


\begin{thebibliography}{20}

\bibitem{AK} F. Albiac and N. Kalton.
Topics in Banach space theory. Second edition. % With a foreword by Gilles Godefroy . Grad. Texts in Math., 233
Springer, [Cham], 2016. % xx+508 pp.

\bibitem{AGM} S. Argyros, M. Gonzalez, and P. Motakis.
Embedding $\ell_2$ and $J$ into subspaces of $JT$ and $JT^*$. Preprint, \url{https://arxiv.org/abs/2603.17886}

%\bibitem{AuSz} G. Aubrun and S. Szarek.Alice and Bob meet Banach. The interface of asymptotic geometric analysis and quantum information theory. % Math. Surveys Monogr., 223
%American Mathematical Society, Providence, RI, 2017. % xxi+414 pp.

\bibitem{ABT} Y. Azouzi, A. Ben Rjeb, and P. Tradacete.
The strong Nakano property in Banach lattices.
Rev. Real Acad. Cienc. Exactas Fis. Nat. Ser. A -- Mat. 119, 99 (2025).
% \url{https://arxiv.org/abs/2401.15629}

\bibitem{AD} Y. Azouzi and W. Dhifaoui.
Order denseness in free Banach lattices.
\url{https://arxiv.org/abs/2508.11648}

\bibitem{BDGJN}
G. Bennett, L. Dor, V. Goodman, W. Johnson, and C. Newman.
On uncomplemented subspaces of $L_p$, $1<p<2$.
Israel J. Math. 26 (1977), 178--187.

\bibitem{BL} Y. Benyamini and J. Lindenstrauss.
Geometric nonlinear functional analysis. Vol. 1. % Amer. Math. Soc. Colloq. Publ., 48
American Mathematical Society, Providence, RI, 2000. % xii+488 pp.

\bibitem{bourgain-DP} J. Bourgain. 
Dunford-Pettis operators on $L^1$ and the Radon-Nikodym property.
Israel J. Math. 37 (1980) 34--47.

\bibitem{Bour} R. Bourgin. Geometric aspects of convex sets with the Radon-Nikod\'ym property.
% Lecture Notes in Math., 993
Springer-Verlag, Berlin, 1983. % xii+ 474 pp.

\bibitem{Car} C.S. Cardassi.
Strictly $p$-integral and $p$-nuclear operators. Analyse harmonique: Groupe de Travail sur les Espaces de Banach Invariants par Translation, Exp. No. 2, 22 pp.
Publ. Math. Orsay, 89-02 [Mathematical Publications of Orsay]
Universit\'e de Paris-Sud, D{\'e}partement de Mathématiques, Orsay, 1989.

\bibitem{DLOT} G. Dales, N. Laustsen, T. Oikhberg, and V. Troitsky.
Multi-norms and Banach lattices.
Dissertationes Math.   524 (2017), 115 pp.

\bibitem{dPWi} B. de Pagter and A. Wickstead.
Free and projective Banach lattices.
Proc. Roy. Soc. Edinburgh Sect. A 145 (2015), 105--143.

\bibitem{DeFl} A. Defant and K. Floret.
Tensor norms and operator ideals.  % North-Holland Math. Stud., 176
North-Holland Publishing Co., Amsterdam, 1993. % xii+566 pp.

\bibitem{DJT}
J.~Diestel, H.~Jarchow, and A.~Tonge, \emph{Absolutely summing operators},
% Cambridge Studies in Advanced Mathematics, vol.~43,
Cambridge University Press, Cambridge, 1995. % \MR{1342297}.

\bibitem{DU-survey} J. Diestel and J. Uhl. 
The Radon-Nikodym theorem for Banach space valued measures.
Rocky Mountain J. Math.   6 (1976), no. 1, 1–46.

\bibitem{DU} J. Diestel and J. Uhl. Vector measures.%  Math. Surveys, No. 15
American Mathematical Society, Providence, RI, 1977. % xiii+322 pp.

% \bibitem{DGKR} S. Dilworth, C. Gartland, D. Kutzarova, and N. Randrianarivony. Nondentable sets in Banach spaces. J. Convex Anal. 28 (2021), 31–40.

\bibitem{FLM} T. Figiel, J. Lindenstrauss, and V. Milman.
The dimension of almost spherical sections of convex bodies.
Acta Math. 139 (1977), 53--94.

\bibitem{Gar} D.J.H. Garling.
Diagonal mappings between sequence spaces.
Studia Math. 51 (1974), 129--138.

\bibitem{GSLTT} E. Garcia-Sanchez,D. Leung, M. Taylor, and P. Tradacete.
Banach lattices with upper $p$-estimates: free and injective objects.
Math. Ann. 391 (2025), 3363--3398.


\bibitem{GS} N. Ghoussoub, and E. Saab. On the weak Radon-Nikod\'ym property. Proc. Amer. Math. Soc. 81.1 (1981), 81-84.

\bibitem{GTT} E. Gluskin, N. Tomczak-Jaegermann, and L. Tzafriri.
Subspaces of $\ell^N_p$ of small codimension.
Israel J. Math. 79 (1992), 173--192.

% \bibitem{GKS} G. Godefroy, N. Kalton, and P. Saphar. Unconditional ideals in Banach spaces. Studia Math. 104, 13--59 (1993).

\bibitem{Hag73} J. Hagler. Some more Banach spaces which contain $L^1$. Studia Math. 46 (1973), 35--42.

\bibitem{Hein} S. Heinrich. Ultraproducts in Banach space theory.
J. Reine Angew. Math. 313 (1980), 72--104.

%\bibitem{James-bush} R.C.~James.
%Subbushes and extreme points in Banach spaces.
%Proceedings of research workshop on Banach space theory (Iowa City, Iowa, 1981), 59--81.
%University of Iowa, Iowa City, IA, 1982

\bibitem{JLTTT} 
H.~Jard\'on-S\'anchez, N.J.~Laustsen, M.A.~Taylor, P.~Tradacete, and V.G.~Troitsky, \emph{Free Banach lattices under convexity conditions},  Rev. R. Acad. Cienc. Exactas F\'is. Nat. Ser. A Mat. RACSAM 116 (2022), no. 1, Paper No. 15.

%\bibitem{Ka84} N. Kalton. Locally complemented subspaces and ${\mathcal{L}}_p$-spaces for $0 < p < 1$. Math. Nachr.   115 (1984), 71--97.

\bibitem{LaTr} N.J. Laustsen and P. Tradacete. Banach lattices of positively homogeneous functions induced by a Banach space, Proc. Edinb. Math. Soc., 69(2) (2026), 544–580.

\bibitem{LiRo} J. Lindenstrauss and H. Rosenthal. The ${\mathcal L}_p$ spaces.
Israel J. Math. 7 (1969), 325--349.


\bibitem{LiSt} J. Lindenstrauss and C. Stegall. Examples of separable spaces which do not contain $\ell_1$ and whose duals are non-separable.
Studia Math.   54 (1975), 81--105.

\bibitem{LT1} J. Lindenstrauss and L. Tzafriri, Classical Banach Spaces, Vol. I, Springer Verlag, 1977.

% \bibitem{LT2} Y.~Lindenstrauss and L. Tzafriri. vol II

%\bibitem{LoR} H. Lotz and H. Rosenthal. Embeddings of $C(\Delta )$ and $L\sp{1}[0,\,1]$ in Banach lattices. Israel J. Math. 31 (1978), 169--179.

\bibitem{Maurey} B. Maurey. 
Type, cotype and K-convexity. Handbook of the geometry of Banach spaces, Vol. 2, 1299--1332.
North-Holland, Amsterdam, 2003.

\bibitem{M-N} P. Meyer-Nieberg. Banach lattices. Springer-Verlag, Berlin, 1991.

\bibitem{MS} V. Milman and G. Schechtman.
Asymptotic theory of finite-dimensional normed spaces. 
% With an appendix by M. Gromov. Lecture Notes in Mathematics, 1200.
Springer-Verlag, Berlin, 1986. 

\bibitem{Musial} K. Musia\l{}. The weak Radon-Nikodym property in Banach spaces. Studia Math. 64 (1979), 151--174.

\bibitem{OiGEO} T. Oikhberg. 
Geometry of unit balls of free Banach lattices, and its applications.
J. Funct. Anal. 286 (2024), % no.8, 
Paper No. 110351, 27 pp.

\bibitem{OiTTT} T. Oikhberg, M. Taylor, P. Tradacete, and V. Troitsky. Free Banach lattices. J. Eur. Math. Soc. (JEMS) 28 (2026), no. 10, 4387--4514.

\bibitem{OiTTT-cor} T. Oikhberg, M. Taylor, P. Tradacete, and V. Troitsky. Corrigendum to ``Free Banach lattices.'' J. Eur. Math. Soc. (JEMS) 28 (2026), no. 10, 4515--4518.

\bibitem{OiTu} T.~Oikhberg and M.A.~Tursi.
Order extreme points and solid convex hulls, The mathematical legacy of Victor Lomonosov -- operator theory, 297--315, Adv. Anal. Geom., 2, De Gruyter, Berlin, 2020. 
% Renorming AM-spaces, Proceedings of the AMS, 150 (2022), no. 3, 1127--1139.

\bibitem{OjP} E. Oja and M. Poldvere.
Principle of local reflexivity revisited.
Proc. Amer. Math. Soc. 135 (2007), 1081--1088. 

% \bibitem{Per} A. Persson. On some properties of $p$-nuclear and $p$-integral operators. Studia Math. 33 (1969), 213--222.

\bibitem{Phelps} R.~Phelps.
Lectures on Choquet's theorem. Second edition.
% Lecture Notes in Math., 1757. 
Springer, Berlin, 2001. % viii+124 pp.

%\bibitem{PieID} A. Pietsch. Operator ideals.
% Translated from German by the author North-Holland Math. Library, 20
%North-Holland Publishing Co., Amsterdam-New York, 1980. % 451 pp.

% \bibitem{Pisier-vol} G. Pisier. The volume of convex bodies and Banach space geometry.
% % Cambridge Tracts in Math., 94
% Cambridge University Press, Cambridge, 1989. xvi+250 pp.

%\bibitem{Pisier-interp} G. Pisier.
%Complex interpolation between Hilbert, Banach and operator spaces. % (English summary)
%Mem. Amer. Math. Soc.   208 (2010), no. 978, vi+78 pp.

%\bibitem{RU} L. H. Riddle, and J. J. Uhl Jr. Martingales and the fine line between Asplund spaces and spaces not containing a copy of $\ell_1$. Martingale Theory in Harmonic Analysis and Banach Spaces: Proceedings of the NSF-CBMS Conference Held at the Cleveland State University, Cleveland, Ohio, July 13–17, 1981. Springer Berlin Heidelberg, 1982.

% \bibitem{Stegall} C. Stegall. The Radon-Nikod\'ym property in Banach spaces. I.
% Proceedings of the conferences on vector measures and integral representations of operators, and on functional analysis/Banach space geometry (Essen, 1982), 1--61.
% Vorlesungen Fachbereich Math. Univ. Essen, 10 [Lecture Notes in Mathematics at the University of Essen], Universit\"at Essen, Fachbereich Mathematik, Essen, 1983

% \bibitem{SteRe} C. Stegall and J. Retherford. Fully nuclear and completely nuclear operators with applications to ${\mathcal L}_1$ and ${\mathcal L}_\infty$-spaces. Trans. Amer. Math. Soc. 163 (1972), 457--492.

% \bibitem{SiYo} B. Sims and D. Yost. Linear Hahn-Banach extension operators. Proc. Edinburgh Math. Soc. (2)   32 (1989), 53--57.

\bibitem{Talagrand} M. Talagrand. Pettis integral and measure theory. Mem. Amer. Math. Soc. Vol. 307. Providence, RI: American Mathematical Society, 1984.

\bibitem{Taylor} M. Taylor.
Unbounded convergences in vector lattices. MSc thesis, University of Alberta, 2018.

\bibitem{Wick07} A. Wickstead.
An isomorphic version of Nakano's characterisation of $C_0(\Sigma)$.
Positivity  11 (2007), 609--615.


\end{thebibliography}


\end{document}


As an application, we consider the evaluation of the $\pi_p(\id_p)$, where $\id_p$ is the formal identity $\ell_p \to c_0$. Pietsch Factorization gives $\pi_p(\id_p) \leq 1$, and conversely, $1 = \|\id_p\| \leq \pi_p(\id_p)$ (the norm is witnessed on vectors with singleton support, where the $\infty$ and $p$ norms are equal). Below we show that such vectors are essential for estimating $\pi_p(\id_p)$.

\begin{proposition}
	For every $p \in (1,\infty)$, and every $\vr > 0$ there exists $\delta > 0$, with the following property: if $z_1, \ldots, z_N \in \ell_p$ satisfy $\|(z_i)\|_{p, {\textrm{weak}}} \leq 1$ and $\sup_i \|z_i\|_\infty/\|z_i\|_p \leq \delta$, then $\sum_i \|z_i\|_\infty^p \leq \vr^p$. 
\end{proposition}

Here, $\| \cdot \|_q$ stands for the $\ell_q$-norm, and as before, $\|(z_i)\|_{p, {\textrm{weak}}} = \sup \big\{ \sum_i |\langle z^*, z_i\rangle|^p : \|z^*\|_{p'} \leq 1 \big\} \leq 1$, $p' = p/(p-1)$.

\begin{proof}
	Suppose, for the sake of contradiction, that there exists $\vr > 0$ so that, for any $m \in \N$, there exist $z_{1,m}, \ldots, z_{N_m,m} \in \ell_p^{K_m}$ so that $\| (z_{i,m})_{i=1}^{N_m} \|_{{p,\textrm{weak}}} = 1$, and, for any $i$, $\|z_{i,m}\|_\infty < 2^{-m} \|z_{i,m}\|_p$.
	Identify $\ell_p$ with $(\oplus_m \ell_p^{K_m})_p$, and denote the canonical projection onto the $m$-th summand $\ell_p^{K_m}$ by $P_m$. For $z = (z_m) \in (\oplus_m \ell_p^{K_m})_p$, define
	$$
	f(z) = \min\{ \|z\|_\infty, \max_m 2^{-m} \|z_m\| \} ,
	$$
	and note that $f(z_{i,m}) = \|z_{i,m}\|_\infty$ for any $m$ and any $i \in \{1, \ldots, N_m\}$.
	
	Let $X = \ell_{p'}$, then $\ell_p = X^*$. It is clear that $f$ is weakly (equivalently, weak$^*$) continuous on bounded sets (on such sets, the weak convergence is metrizable, and coincides with componentwise convergence). Moreover, $f \in H_{w^*}^p[X]$, since $0 \leq f(z) \leq \|\id_p z\|$; this shows that $\|f\|_p \leq 1$. 
	
	Now consider $g \in I_{w^*}[X]$. By \Cref{good approx}, we can assume that $g$ is a lattice expression in $(y_\ell){\ell=1}^L$, where $P_m y_\ell = 0$ for any $\ell$, and any $m > m_0$. Fixing such $m$, we see that
	$$
	\|f-g\|_p^p \geq \sum_i \big| f(z_{i,m}) - g(z_{i,m}) \big|^p = \vr^p .
	$$
	Thus, $f$ cannot be approximated by the elements of $I_{w^*}[X]$ in the $p$-norm, which contradicts \Cref{la-tr}.
\end{proof}


======================================================


\bibitem{AuSz} G. Aubrun and S. Szarek.Alice and Bob meet Banach. The interface of asymptotic geometric analysis and quantum information theory. % Math. Surveys Monogr., 223
American Mathematical Society, Providence, RI, 2017. % xxi+414 pp.

\bibitem{Beau} B. Beauzamy.
Introduction to Banach spaces and their geometry.
% Notas Mat.[Mathematical Notes] North-Holland Math. Stud., 68
North-Holland Publishing Co., Amsterdam-New York, 1982. % xi+308 pp.

\bibitem{Cas} P. Casazza.
Approximation properties.
Handbook of the geometry of Banach spaces, Vol. I, 271--316.
North-Holland Publishing Co., Amsterdam, 2001

\bibitem{CK-AP} P. Casazza and N. Kalton. 
Notes on approximation properties in separable Banach spaces.
Geometry of Banach spaces (Strobl, 1989), 49--63. % London Math. Soc. Lecture Note Ser., 158
Cambridge University Press, Cambridge, 1990

\bibitem{DFS} J. Diestel, J. Fourie, and J. Swart.
The metric theory of tensor products. Grothendieck's r\'esum\'e revisited.
American Mathematical Society, Providence, RI, 2008.

\bibitem{Dil} S. Dilworth. 
Special Banach lattices and their applications.
Handbook of the geometry of Banach spaces, Vol. I, 497--532.
North-Holland Publishing Co., Amsterdam, 2001

\bibitem{FHHMZ11}
M. Fabian, P. Habala, P. Hajek, Petr, V. Montesinos, and V. Zizler.
Banach space theory. The basis for linear and nonlinear analysis.
% CMS Books Math./Ouvrages Math. SMC
Springer, New York, 2011. % xiv+820 pp.

\bibitem{ER} E. Effros and Z.-J. Ruan.
Operator spaces. 
London Mathematical Society Monographs.
% New Series, 23. The Clarendon Press, 
Oxford University Press, New York, 2000.

\bibitem{Enc} 
Encyclopaedia of mathematics. Supplement. Vol. I. 
Edited by M. Hazewinkel. Kluwer Academic Publishers, Dordrecht, 1997.

\bibitem{Joh} K. John.On the compact non-nuclear operator problem. Math. Ann. 287 (1990), 509--514. 

\bibitem{JoDEC} W. Johnson.
Finite-dimensional Schauder decompositions in $\pi_\lambda$ and dual $\pi_\lambda$ spaces.
Illinois J. Math. 14 (1970), 642--647.

\bibitem{JoFAC} W. Johnson.
Factoring compact operators. Israel J. Math. 9 (1971), 337--345.

\bibitem{Lin-nonsep-refl} Y.~Lindenstrauss. On nonseparable reflexive Banach spaces.
Bull. Amer. Math. Soc. 72 (1966), 967--970. 

\bibitem{LT1} Y.~Lindenstrauss and L. Tzafriri. vol I

\bibitem{LT2} Y.~Lindenstrauss and L. Tzafriri. vol II

\bibitem{MS} V. Milman and G. Schechtman.
Asymptotic theory of finite-dimensional normed spaces. 
% With an appendix by M. Gromov. Lecture Notes in Mathematics, 1200.
Springer-Verlag, Berlin, 1986. 

\bibitem{Mus} J. Musielak.
Orlicz spaces and modular spaces. % Lecture Notes in Math., 1034
Springer-Verlag, Berlin, 1983. % iii+222 pp.


\bibitem{OZ} E. Oja and I. Zolk. 
The asymptotically commuting bounded approximation property of Banach spaces.
J. Funct. Anal. 266 (2014), 1068--1087.

\bibitem{PisFIN} G. Pisier.
Un th\'eor\`eme sur les op\'erateurs lin\'eaires entre espaces de Banach
 qui se factorisent par un espace de Hilbert.
Ann. Sci. \'Ecole Norm. Sup. (4) 13 (1980), 23--43.

\bibitem{PisINF} G. Pisier.
Counterexamples to a conjecture of Grothendieck. 
Acta Math. 151 (1983), 181--208. 

\bibitem{Rud} M. Rudelson.
Estimates on the weak distance between finite dimensional Banach spaces.
Israel J. Math. 89 (1995), 189--204.

\bibitem{Rya} R. Ryan.
Introduction to tensor products of Banach spaces.
% Springer Monographs in Mathematics. 
Springer-Verlag, London, 2002.

\bibitem{Str} W. Stromquist.
The maximum distance between two-dimensional Banach spaces. 
Math. Scand. 48 (1981), 205--225. 

\bibitem{T-J} N. Tomczak-Jaegermann.
Banach-Mazur distances and finite-dimensional operator ideals.
Longman Scientific and Technical, 1989.
